\section{Exponential–family conjugate single–segment evidence and moments}
\label{sec:ef}

\paragraph{Plate model (single segment).}
For a block $(i,j]$ and conjugate hyperparameters $(\alpha_0,\beta_0)$, the generative model is
\begin{equation}
\theta \sim p(\theta\mid \alpha_0,\beta_0),
\qquad
y_t\mid \theta \stackrel{\mathrm{ind}}{\sim} p(y_t\mid\theta,w_t) \text{ as in \eqref{eq:EFweighted}},\quad t\in\{i+1,\dots,j\}.
\label{eq:plate-ef}
\end{equation}
The block evidence $A^{(0)}_{ij}$ is obtained by marginalizing $\theta$ out of \eqref{eq:plate-ef}.
\begin{figure}[H]
\centering
\begin{tikzpicture}
  \node[obs] (y) {$y_t$};
  \node[latent, above=0.9cm of y] (theta) {$\theta$};
  \node[const, left=0.9cm of theta] (ab) {$\alpha_0,\beta_0$};
  \node[const, right=0.9cm of y] (w) {$w_t$};
  \edge {theta} {y};
  \edge {ab} {theta};
  \edge {w} {y};
  \plate {plate1} {(y)} {$t\in(i,j]$};
\end{tikzpicture}
\caption[Single conjugate EF block plate model]{Plate model for a single conjugate EF block. Shaded nodes are observed; unshaded are latent; point nodes are fixed hyperparameters/exposures. Marginalizing $\theta$ gives $A^{(0)}_{ij}$.}
\label{fig:plate-ef}
\end{figure}

Recall the generic EF--conjugate specification \citep{diaconis1979conjugate,bernardo1994bayesian}. Conditional on a (possibly low-dimensional)
parameter $\theta$, observations $y_t$ in a block $(i,j]$ are independent with density
\begin{equation}
p(y_t\mid\theta)
  = \exp\{w_t[T(y_t)^\top\eta(\theta) - A(\theta)] + \log h(y_t,w_t)\},
\label{eq:EFweighted}
\end{equation}
where $w_t\ge 0$ are exposure weights and $h$ collects base measures. The conjugate prior for
$\theta$ is
\[
p(\theta\mid\alpha,\beta)
  = \exp\{\eta(\theta)^\top\alpha - \beta A(\theta) - \log Z(\alpha,\beta)\},
\]
with normalizing integral $Z(\alpha,\beta)$ defined by \eqref{eq:Zconvention}.

For a block $(i,j]$ define the cumulative sums
\[
S_{ij} := \sum_{t=i+1}^j w_t\,T(y_t),\qquad
W_{ij} := \sum_{t=i+1}^j w_t,\qquad
H_{ij} := \sum_{t=i+1}^j \log h(y_t,w_t).
\]

\begin{theorem}[EF--conjugate block integration]
\label{thm:ef-integral}
Fix a block $(i,j]$ and a conjugate prior $p(\theta\mid\alpha_0,\beta_0)$. Assume the integral
$Z(\alpha_0+S_{ij},\beta_0+W_{ij})$ is finite for every $(i,j]$ with $0\le i<j\le n$, and that the
moment $\mathbb{E}[m(\theta)^r\mid\alpha_0+S_{ij},\beta_0+W_{ij}]$ is finite for each required
$r\ge 1$. Under the weighted EF model \eqref{eq:EFweighted}, the integrated single--segment
evidence and the $r$th moment integral of the observation-scale EF mean
$m(\theta):=\mathbb{E}[Y\mid\theta]$ are
\begin{align}
A^{(0)}_{ij}
 &= \int \Big[\prod_{t=i+1}^j p(y_t\mid\theta)\Big]
      p(\theta\mid\alpha_0,\beta_0)\,d\theta \nonumber\\
 &= \exp\{H_{ij}\}\,
      \frac{Z(\alpha_0 + S_{ij},\beta_0 + W_{ij})}{Z(\alpha_0,\beta_0)},
      \label{eq:A0}\\
A^{(r)}_{ij}
 &= \int \Big[\prod_{t=i+1}^j p(y_t\mid\theta)\Big]
      p(\theta\mid\alpha_0,\beta_0)\,m(\theta)^r\,d\theta \nonumber\\
 &= A^{(0)}_{ij}\,
    \mathbb{E}\!\big[m(\theta)^r
      \,\big|\,\alpha_0+S_{ij},\beta_0+W_{ij}\big],
      \qquad r=1,2,\dots.
      \label{eq:Ar}
\end{align}
\end{theorem}

\begin{proof}
\ProofStep{Step 1: factorize the block likelihood.}
By independence and \eqref{eq:EFweighted},
\[
\prod_{t=i+1}^j p(y_t\mid\theta)
  = \exp\!\Big\{
      \sum_{t=i+1}^j w_t[T(y_t)^\top\eta(\theta) - A(\theta)]
      + \sum_{t=i+1}^j \log h(y_t,w_t)
     \Big\}.
\]
Collecting block sums gives
\[
\prod_{t=i+1}^j p(y_t\mid\theta)
  = \exp\{H_{ij} + \eta(\theta)^\top S_{ij} - W_{ij} A(\theta)\}.
\]

\ProofStep{Step 2: multiply by the conjugate prior and identify the posterior kernel.}
Multiply by $p(\theta\mid\alpha_0,\beta_0)=\exp\{\eta(\theta)^\top\alpha_0 - \beta_0 A(\theta) - \log Z(\alpha_0,\beta_0)\}$:
\[
\Big[\prod_{t=i+1}^j p(y_t\mid\theta)\Big] p(\theta\mid\alpha_0,\beta_0)
 = \frac{\exp\{H_{ij}\}}{Z(\alpha_0,\beta_0)}\,
   \exp\{\eta(\theta)^\top(\alpha_0+S_{ij}) - (\beta_0+W_{ij})A(\theta)\}.
\]

\ProofStep{Step 3: integrate by using the normalizer definition.}
By the definition \eqref{eq:Zconvention},
\[
\int \exp\{\eta(\theta)^\top(\alpha_0+S_{ij}) - (\beta_0+W_{ij})A(\theta)\}\,d\theta
  = Z(\alpha_0+S_{ij},\beta_0+W_{ij}).
\]
Therefore,
\[
A^{(0)}_{ij}
 = \exp\{H_{ij}\}\,\frac{Z(\alpha_0+S_{ij},\beta_0+W_{ij})}{Z(\alpha_0,\beta_0)},
\]
which is \eqref{eq:A0}.

\ProofStep{Step 4: moment integrals.}
For $r\ge 1$, the integrand defining $A^{(r)}_{ij}$ equals the integrand of $A^{(0)}_{ij}$
times $m(\theta)^r$. Dividing by $A^{(0)}_{ij}$ therefore yields the expectation of
$m(\theta)^r$ under the posterior
$p(\theta\mid \alpha_0+S_{ij},\beta_0+W_{ij})$, which is finite by the hypothesis of the theorem,
giving \eqref{eq:Ar}.
\end{proof}

\section{Dynamic programming: evidence, boundaries, and moments}
\label{sec:dp}

For $k\in\{0,\dots,k_{\max}\}$ and $j\in\{0,\dots,n\}$, let $L_{k j}$ denote the evidence of
the prefix $y_{1:j}$ under a segmentation into $k$ segments (after integrating out segment
parameters), and let $R_{k i}$ denote the evidence of the suffix $y_{i+1:n}$ under $k$
segments. We adopt the conventions
\[
L_{0,0}=1,\quad L_{0,j}=0\ (j>0),\qquad
R_{0,n}=1,\quad R_{0,i}=0\ (i<n).
\]

\subsection{Forward/backward recursions}
Given single-block evidences $A^{(0)}_{ij}$, the forward and backward recursions are
\begin{equation}
L_{k+1,j}
 = \sum_{h=k}^{j-1} L_{k,h}\,A^{(0)}_{h j},\qquad
R_{k+1,i}
 = \sum_{h=i+1}^{n-k} A^{(0)}_{i h}\,R_{k,h},
\label{eq:LR}
\end{equation}
for $k\ge0$, $1\le j\le n$, and $0\le i\le n-1$, where sums are restricted to valid indices.


\paragraph{Implementation.}
Algorithm~\ref{alg:dp-core} presents the core forward/backward recursions in a numerically stable log-space form.
The routine takes as input a triangular array of block log-evidences $\log \mathcal{L}_{ij}$ (from Section~\ref{sec:ef} or Sections~\ref{sec:families}--\ref{sec:nonconj}) and returns the forward messages, backward messages, and segment-count posteriors needed for boundary and moment extraction.

\begin{algorithm}[H]
\caption{Generic DP for EF segmentation (prefix/suffix form)}
\label{alg:dp-core}
\KwIn{Block evidences $\widetilde{A}^{(0)}_{ij}$ for $0\le i<j\le n$; $k_{\max}$; $C_k$ (default $C_k=\binom{n-1}{k-1}$)}
Initialize $\widetilde{L}_{0,0}=1$, $\widetilde{L}_{0,j>0}=0$; $\widetilde{R}_{0,n}=1$, $\widetilde{R}_{0,i<n}=0$\;
\For{$k=0$ \KwTo $k_{\max}-1$}{
  \For{$j=1$ \KwTo $n$}{
    $\widetilde{L}_{k+1,j}\leftarrow \sum_{h=k}^{j-1} \widetilde{L}_{k,h}\widetilde{A}^{(0)}_{h j}$\;
  }
  \For{$i=0$ \KwTo $n-1$}{
    $\widetilde{R}_{k+1,i}\leftarrow \sum_{h=i+1}^{n-k} \widetilde{A}^{(0)}_{i h} \widetilde{R}_{k,h}$\;
  }
}
Compute $P(k\mid y)\propto p(k)\,\widetilde{L}_{k n}/C_k$ (cf.\ \eqref{eq:post-k}) and choose $\widehat{k}$\;
\For{$p=1$ \KwTo $\widehat{k}-1$}{
  For $h\in\{p,\dots,n-(\widehat{k}-p)\}$ set $P(t_p=h\mid y,\widehat{k})\leftarrow \widetilde{L}_{p,h}\widetilde{R}_{\widehat{k}-p,h}/\widetilde{L}_{\widehat{k}n}$\;
  Record the marginal mode $\arg\max_h P(t_p=h\mid y,\widehat{k})$ as a diagnostic only\;
}
\KwOut{$\{\widetilde{L}_{k j}\}$, $\{\widetilde{R}_{k i}\}$, $P(k\mid y)$, $\widehat{k}$, and boundary-marginal diagnostics}
\end{algorithm}


\subsection{Incorporating segment-factorized boundary priors}
\label{subsec:priorinc}

Suppose the boundary prior given $k$ factorizes over segment lengths as in \eqref{eq:lengthprior}:
\[
p(t\mid k) = \frac{1}{C_k}\prod_{q=1}^k g(\Delta_x(t_{q-1},t_q)),
\qquad
C_k := \sum_{t:\,0=t_0<\cdots<t_k=n}\ \prod_{q=1}^k g(\Delta_x(t_{q-1},t_q)).
\]
Then we can absorb $g$ into block evidences by defining
\begin{equation}
\widetilde{A}^{(0)}_{ij} := A^{(0)}_{ij}\,g(\Delta_x(i,j)),
\qquad
\widetilde{A}^{(r)}_{ij} := A^{(r)}_{ij}\,g(\Delta_x(i,j)).
\label{eq:Atilde}
\end{equation}
Running the DP \eqref{eq:LR} with $\widetilde{A}^{(0)}_{ij}$ computes $\widetilde{L}_{k n}$,
the \emph{unnormalized} evidence sum over all segmentations weighted by $g$. The normalizer
$C_k$ can be computed by the same DP with $A^{(0)}_{ij}\equiv 1$:
\begin{equation}
C_k
 = \sum_{t}\prod_{q=1}^k g(\Delta_x(t_{q-1},t_q))
 = L^{(g)}_{k n},
\quad\text{where}\quad
L^{(g)}_{k+1,j}=\sum_{h=k}^{j-1}L^{(g)}_{k,h}\,g(\Delta_x(h,j)),\ L^{(g)}_{0,0}=1.
\label{eq:Ck-general}
\end{equation}
In the index-uniform case $g\equiv 1$, \eqref{eq:Ck-general} reduces to $C_k=\binom{n-1}{k-1}$.

\subsection[Posteriors over segment counts, boundary marginals, and joint MAP segmentations]{Posteriors over $k$, boundary marginals, and joint MAP segmentations}
Let $p(k)$ be a prior on segment count. Under index-uniform $p(t\mid k)$ we have
$C_k=\binom{n-1}{k-1}$ and
\[
P(y\mid k) = \frac{L_{k n}}{C_k},
\qquad
P(k\mid y)\propto p(k)\,\frac{L_{k n}}{C_k}.
\]
For general segment-factorized $p(t\mid k)$, replace $L$ by $\widetilde{L}$ and $C_k$ by
\eqref{eq:Ck-general}. In all cases we can write compactly
\begin{equation}
P(k\mid y)\ \propto\ p(k)\,\frac{\widetilde{L}_{k n}}{C_k},
\qquad k=1,\dots,k_{\max}.
\label{eq:post-k}
\end{equation}
Let $\widehat{k}=\arg\max_k P(k\mid y)$ denote the MAP segment count.

Conditional on $k$, the marginal posterior for the $p$th boundary $t_p$ is
\begin{equation}
P(t_p=h\mid y,k)
 = \frac{\widetilde{L}_{p,h}\,\widetilde{R}_{k-p,h}}{\widetilde{L}_{k n}},
\qquad
h\in\{p,\dots,n-(k-p)\},
\label{eq:boundary-post}
\end{equation}
where $(\widetilde{L},\widetilde{R})$ are computed from $\widetilde{A}^{(0)}$ as in
\eqref{eq:LR}. The denominator is the unnormalized evidence under the length-weighted partition
prior; the normalizer $C_k$ cancels in \eqref{eq:boundary-post} because the same $C_k$ multiplies
every boundary vector in the numerator and denominator. (When $g\equiv 1$, tildes can be dropped.)

\begin{remark}[Boundary-event marginal]
\label{rem:boundary-event}
A related object we use repeatedly is the \emph{boundary-event marginal}
\begin{equation}
P(b_i=1\mid y,k)
 := \sum_{p=1}^{k-1} P(t_p=i\mid y,k)
 = \sum_{p=1}^{k-1}\frac{\widetilde{L}_{p,i}\,\widetilde{R}_{k-p,i}}{\widetilde{L}_{k n}},
 \qquad 1\le i\le n-1,
\label{eq:boundary-event}
\end{equation}
the posterior probability that interior index $i$ coincides with \emph{some} boundary, without
conditioning on which boundary. The objects $P(t_p=h\mid y,k)$ and $P(b_i=1\mid y,k)$ are
distinct: the former is a probability distribution over $h$ for each fixed $p$ (and sums to $1$ over
$h$), while the latter is a per-index Bernoulli probability summing to $k-1$ over $i$. Calibration
experiments in Section~\ref{sec:experiments} target $P(b_i=1\mid y,k)$; a full derivation is given
in Appendix~\ref{app:marginals}.
\end{remark}

A useful diagnostic is the \emph{marginal boundary mode}
\[
\widetilde{t}^{\,\mathrm{marg}}_p(k)
  \in \arg\max_h P(t_p=h\mid y,k).
\]
However, maximizing the marginals independently need not yield an ordered boundary vector and, in
general, does \emph{not} recover the joint MAP segmentation.

\begin{remark}[A small counterexample]
\label{rem:marg-vs-joint}
Consider $n=5$ and $k=3$, so the two interior boundaries satisfy $1\le t_1<t_2\le4$.
For example, let the joint posterior put masses $0.30,0.28,0.22,0.20$ on the four ordered
configurations $(1,4),(2,3),(2,4),(3,4)$, respectively. The marginal modes are then
$t_1=2$ and $t_2=4$, so the independently selected marginal-mode vector is $(2,4)$, even though
the joint MAP configuration is $(1,4)$. This is exactly the kind of discrepancy the max-sum DP
resolves.
\end{remark}

The joint MAP segmentation is the solution of
\begin{equation}
\widehat{t}^{\,\mathrm{MAP}}(k)
 \in \arg\max_{0=t_0<\cdots<t_k=n}
 \Bigg\{\sum_{q=1}^k \log \widetilde{A}^{(0)}_{t_{q-1} t_q}\Bigg\},
\label{eq:joint-map-k}
\end{equation}
with the convention that any $k$-specific prior terms such as $\log p(k)-\log C_k$ are added only
when selecting among different values of $k$. Equation~\eqref{eq:joint-map-k} is solved exactly by
a max-sum DP with backpointers. Inline, the max-sum recursion reads
\begin{equation}
M_{k+1,j} = \max_{h\in\{k,\dots,j-1\}}\big\{M_{k,h} + \log \widetilde{A}^{(0)}_{hj}\big\},
\qquad
h^\star(k+1,j) \in \arg\max_{h} \big\{M_{k,h} + \log \widetilde{A}^{(0)}_{hj}\big\},
\label{eq:max-sum-inline}
\end{equation}
with initial conditions $M_{0,0}=0$, $M_{0,j}=-\infty$ for $j>0$, and backtracking from
$(k,n)$ as detailed in Section~\ref{sec:algorithms}.
In the sequel, whenever we report an \emph{exported MAP segmentation}, we mean the joint optimizer
$\widehat{t}^{\,\mathrm{MAP}}:=\widehat{t}^{\,\mathrm{MAP}}(\widehat{k})$ rather than a vector of marginal
boundary modes.

\subsection{Segment-level and regression-curve posterior moments}
For segment-wise posterior moments attached to a \emph{fixed} segmentation
$t=(t_0,\dots,t_k)$, let segment $q\in\{1,\dots,k\}$ correspond to block
$(i,j]=(t_{q-1},t_q]$. Then
\begin{equation}
\mu^{(r)}_q(t)
 = \frac{\widetilde{A}^{(r)}_{ij}}{\widetilde{A}^{(0)}_{ij}}
 = \frac{A^{(r)}_{ij}}{A^{(0)}_{ij}},
\qquad r=1,2,\dots,
\label{eq:segmom-seg}
\end{equation}
where the final equality holds because the prior factor $g(\Delta_x(i,j))$ cancels. In particular, for
an exported joint MAP segmentation $\widehat{t}^{\,\mathrm{MAP}}$, the segment summaries are obtained by
setting $t=\widehat{t}^{\,\mathrm{MAP}}$ in \eqref{eq:segmom-seg}.

To compute the Bayes regression curve within a fixed $k$ (typically $k=\widehat{k}$),
define, for each block $(i,j]$ and moment order $r$,
\[
F^{(r)}_{ij}(k)
 := \frac{1}{\widetilde{L}_{k n}}\sum_{m=1}^k
       \widetilde{L}_{m-1,i}\,\widetilde{A}^{(r)}_{ij}\,\widetilde{R}_{k-m,j}.
\]
Then the $r$th posterior moment of the latent mean at index $t$ given $k$ is obtained by summing
$F^{(r)}_{ij}(k)$ over all candidate blocks $(i,j]$ that contain the index $t$ (i.e., all blocks
with $i<t\le j$), so that each candidate segmentation contributes exactly once through the block
that actually covers $t$:
\begin{equation}
\widehat{\mu}'^{(r)}_t
 := \mathbb{E}[(\mu'_t)^r\mid y,k]
 = \sum_{i<t\le j} F^{(r)}_{ij}(k),
\qquad t=1,\dots,n.
\label{eq:segmom}
\end{equation}
In particular, $\widehat{\mu}'_t := \widehat{\mu}'^{(1)}_t$ is the Bayes regression curve and
$\mathbb{V}[\mu'_t\mid y,k]
 = \widehat{\mu}'^{(2)}_t - (\widehat{\mu}'^{(1)}_t)^2$.
If a fully $k$-marginalized curve is desired, compute the same conditional curve for each admissible
$k$ and average the resulting moments using $P(k\mid y)$ from \eqref{eq:post-k}. The empirical
sections report curves conditional on the selected $\widehat k$ unless explicitly stated otherwise.

\begin{theorem}[Correctness of the DP]
\label{thm:dp-correctness}
Assume finite block evidences on all admissible blocks, a finite segment-count range, and a
segment-factorized conditional partition prior satisfying Assumption~\ref{ass:factorizing-prior}
with the boundary-coordinate convention of Section~\ref{sec:notation}. Under the EF--conjugate
single-segment model of Section~\ref{sec:ef}, or under any other block model whose exact evidence
array is supplied, the recursions \eqref{eq:LR} compute exact marginal evidences for all
prefixes/suffixes:
$\widetilde{L}_{k j}$ and $\widetilde{R}_{k i}$ are exact evidence sums over all
$k$-segmentations of the corresponding prefix/suffix. Moreover:
\begin{enumerate}
\item The posterior $P(k\mid y)$ given by \eqref{eq:post-k} is exact.
\item The boundary marginal $P(t_p=h\mid y,k)$ given by \eqref{eq:boundary-post} is exact.
\item For any fixed $k$, the joint MAP segmentation \eqref{eq:joint-map-k} is obtained exactly by the
max-sum recursion $M_{k+1,j}=\max_h\{M_{k,h}+\log\widetilde A^{(0)}_{hj}\}$ with backpointers
$h^\star(k+1,j)$ and subsequent backtracking (the full routine of Section~\ref{sec:algorithms};
detailed proof in Appendix~\ref{app:max-sum-proof}).
\item The segment-wise moments \eqref{eq:segmom-seg} and regression-curve moments \eqref{eq:segmom}
are exact posterior moments conditional on $k$.
\end{enumerate}
\end{theorem}

\begin{proof}
We present the proof for general $\widetilde{A}^{(0)}_{ij}$; the untilded case is the
specialization $g\equiv 1$.

\paragraph{Step 1: interpret $\widetilde{L}_{k j}$.}
Fix $k\ge 1$ and $j\ge 1$. Consider all segmentations of the prefix $y_{1:j}$ into $k$ segments,
i.e.\ all boundary vectors $(t_0,\dots,t_k)$ with $t_0=0$ and $t_k=j$. For any such segmentation,
segment parameters are independent a priori and observations are independent across segments
given the segmentation. Integrating out segment parameters yields a product of block evidences:
\[
p(y_{1:j}\mid t,k)=\prod_{q=1}^k \widetilde{A}^{(0)}_{t_{q-1},t_q}.
\]
Summing over all segmentations gives
\[
\widetilde{L}_{k j}
= \sum_{t:\,t_k=j}\ \prod_{q=1}^k \widetilde{A}^{(0)}_{t_{q-1},t_q}.
\]
Thus $\widetilde{L}_{k j}$ is exactly the marginal evidence sum for the prefix under $k$ segments.

\paragraph{Step 2: derive the forward recursion.}
Any segmentation of $y_{1:j}$ into $k{+}1$ segments has a last boundary at some $h\in\{k,\dots,j-1\}$.
Conditioning on this last boundary splits the segmentation into:
(i) a $k$-segmentation of $y_{1:h}$ and (ii) the final block $(h,j]$.
By the factorization above, the evidence contribution for a fixed $h$ equals
$\widetilde{L}_{k h}\,\widetilde{A}^{(0)}_{h j}$. Summing over all admissible $h$ yields
\[
\widetilde{L}_{k+1,j}=\sum_{h=k}^{j-1} \widetilde{L}_{k h}\,\widetilde{A}^{(0)}_{h j},
\]
which is the left recursion in \eqref{eq:LR}.

\paragraph{Step 3: derive the backward recursion.}
The same argument applies to suffixes by conditioning on the first boundary after $i$, yielding
\[
\widetilde{R}_{k+1,i}=\sum_{h=i+1}^{n-k} \widetilde{A}^{(0)}_{i h}\,\widetilde{R}_{k h},
\]
the right recursion in \eqref{eq:LR}.

\paragraph{Step 4: posterior over $k$.}
For any fixed $k$, the data likelihood marginalized over segmentations is $\widetilde{L}_{k n}$.
If $p(t\mid k)$ includes a normalizer $C_k$ (as in \eqref{eq:lengthprior}), then
$P(y\mid k)=\widetilde{L}_{k n}/C_k$ by definition of $C_k$ as the sum of prior weights across
all segmentations. Bayes' rule with prior $p(k)$ yields \eqref{eq:post-k}.

\paragraph{Step 5: boundary posteriors.}
Fix $k$ and boundary index $p\in\{1,\dots,k-1\}$. Conditioning on $t_p=h$ splits the sequence into:
(i) a prefix $y_{1:h}$ with $p$ segments and (ii) a suffix $y_{h+1:n}$ with $k-p$ segments.
By independence across the split and by the definition of $\widetilde{L},\widetilde{R}$,
the total evidence of all segmentations consistent with $t_p=h$ equals
$\widetilde{L}_{p,h}\,\widetilde{R}_{k-p,h}$. Normalizing by the total evidence
$\widetilde{L}_{k n}$ yields \eqref{eq:boundary-post}.

\paragraph{Step 6: moment formulas.}
For a fixed block $(i,j]$, Theorem~\ref{thm:ef-integral} gives
$A^{(r)}_{ij}=A^{(0)}_{ij}\,\mathbb{E}[m(\theta)^r\mid(i,j]]$. Since the segmentation prior factor
$g(\Delta_x(i,j))$ cancels in the ratio, \eqref{eq:segmom-seg} holds. The regression-curve moment formula
\eqref{eq:segmom} follows by averaging over all segmentations and grouping contributions by the
block $(i,j]$ that contains index $t$.
\end{proof}

\paragraph{Posterior-evidence sanity checks.}
After every DP run, verify four invariants: normalized segment-count weights
sum to one; the full-sequence forward and backward evidences agree, $\widetilde L_{kn}=\widetilde
R_{k0}$; boundary-event marginals satisfy $\sum_{i=1}^{n-1}P(b_i=1\mid y,k)=k-1$; and the
max-sum backtracking score equals the stored terminal max-sum value.

\begin{theorem}[Time and space complexity]
\label{thm:dp-complexity}
Assume single-block evidences $A^{(0)}_{ij}$ are available for all $0\le i<j\le n$.
Then computing $L_{k j}$ and $R_{k i}$ for $k=0,\dots,k_{\max}$ via \eqref{eq:LR} requires
$\mathcal{O}(k_{\max} n^2)$ arithmetic operations. Storing all $A^{(0)}_{ij}$ requires
$\mathcal{O}(n^2)$ memory, and storing the DP arrays requires $\mathcal{O}(k_{\max} n)$.
\end{theorem}

\begin{proof}
For each fixed $k$, the forward recursion computes $L_{k+1,j}$ for $j=1,\dots,n$, and each
$L_{k+1,j}$ sums at most $j-k$ terms. Thus the forward work per $k$ is
$\sum_{j=1}^n \mathcal{O}(j)=\mathcal{O}(n^2)$. The backward recursion is analogous. Repeating
for $k=0,\dots,k_{\max}-1$ gives $\mathcal{O}(k_{\max}n^2)$ time. Memory follows by counting
array sizes.
\end{proof}

\begin{corollary}[Joint MAP is Bayes-optimal under 0--1 loss on the full segmentation]
\label{cor:map-01}
For the loss $\ell_{01}(t,t')=\mathbbm{1}\{t\neq t'\}$, the Bayes estimator $t^\star$ minimizing
$\mathbb{E}[\ell_{01}(t,t')\mid y,k]$ is the joint MAP segmentation
$\widehat{t}^{\,\mathrm{MAP}}(k)$ from \eqref{eq:joint-map-k}. For any additive per-segment loss
$\ell_{\mathrm{add}}(t,t')=\sum_q \ell_q(t_q,t'_q)$ that does not decompose as the pointwise
indicator on the full vector, the Bayes-optimal estimator can differ from
$\widehat{t}^{\,\mathrm{MAP}}(k)$.
\end{corollary}

\begin{proof}
The 0--1 loss's expectation equals $1-p(t=t'\mid y,k)$, which is minimized by the joint posterior
mode, i.e., the joint MAP. Pointwise (say Hamming-type) losses decompose differently and are
minimized by the marginal modes, which we have just noted need not coincide with the joint MAP.
\end{proof}

\begin{remark}[Memory/time trade-offs for large $n$]
\label{rem:mem-time}
Two implementation variants are standard. (i) \emph{Streaming $k$-layers}: keeping only two
consecutive layers of $(L,R)$ gives memory $\mathcal{O}(n)$ (plus $\mathcal{O}(n^2)$ for the
block-evidence matrix) at no time cost, but backpointers and boundary marginals must then be
recomputed. (ii) \emph{Checkpointing}: storing every $\sqrt{k_{\max}}$th $k$-layer gives memory
$\mathcal{O}(\sqrt{k_{\max}}\,n)$ and at most $\mathcal{O}(\sqrt{k_{\max}}\,k_{\max}\,n^2)$ time,
because recomputing from the nearest checkpoint needs at most $\sqrt{k_{\max}}$ forward layers.
When the $n^2$ block-evidence matrix is itself prohibitive, an alternative is to evaluate block
evidences on demand during the DP sweep at cost $\mathcal{O}(k_{\max}n^2)$ total (constant-time
per block via prefix sums).
\end{remark}

\begin{remark}[Evidence across segment counts need not be monotone]
\label{prop:mono-k}
The fixed-$k$ evidence $P(y\mid k)=\widetilde L_{kn}/C_k$ is not monotone in general. Adding a boundary changes both the set of admissible partitions and
the product of block evidences, while the normalized conditional prior contributes the normalizer
$C_k$. Consequently, diagnose segment-count behavior through the normalized evidences
$P(y\mid k)$ and posterior weights $P(k\mid y)$ rather than through an unnormalized pathwise
``add one boundary'' heuristic.
\end{remark}


\subsection{End-to-end BayesBreak inference routine}\label{sec:dp-alg}
Collecting the ingredients from Sections~\ref{sec:ef} and \ref{sec:dp}---block-evidence precomputation,
exact sum-product DP over partitions, and max-sum backtracking for exported segmentations---yields an
end-to-end routine suitable for direct implementation. Algorithm~\ref{alg:bayesbreak} summarizes
this workflow in a modular way: the only model-specific input is a function that evaluates (or
approximates) block evidences $\log \mathcal{L}_{ij}$ and, when needed, block moment numerators.

\begin{algorithm}[H]
\caption{BayesBreak$(y,x,w,\texttt{model},\texttt{hyper},k_{\max},p(k),g)$}
\label{alg:bayesbreak}
\KwIn{
  Observations $y_{1:n}$ at locations $x_{1:n}$ (optional); observation-model weights/exposures $w_{1:n}$ (default determined by the chosen likelihood, often $w_i\equiv 1$);\\
  model family $\texttt{model}\in\{\texttt{Gaussian},\texttt{Poisson},\texttt{Binomial},\texttt{BetaObs},\dots\}$; \\
  hyperparameters $\texttt{hyper}$; maximum segment count $k_{\max}$; prior $p(k)$; length factor $g(\cdot)$ (default $g\equiv 1$).
}
\KwOut{
  Posterior $P(k\mid y)$; MAP segment count $\widehat{k}$; boundary marginals $P(t_p\mid y,\widehat{k})$;\\
  exported joint MAP segmentation $\widehat{t}^{\,\mathrm{MAP}}$; segment-level moments; Bayes regression curve; (optional) evidence $P(y)$.
}
\BlankLine
\textbf{1. Precompute block evidences and moments}\;
\For{$0\le i<j\le n$}{
  Compute block summaries $(S_{ij},W_{ij},H_{ij})$ using prefix sums (or running sums)\;
  Call a family-specific routine to obtain $A^{(0)}_{ij},A^{(1)}_{ij},A^{(2)}_{ij}$\;
  Set $\widetilde{A}^{(r)}_{ij}\leftarrow A^{(r)}_{ij}\,g(\Delta_x(i,j))$ for $r\in\{0,1,2\}$ (cf.\ \eqref{eq:Atilde})\;
}
Compute $C_k$ under the same admissibility and length-prior convention used in the score array; for $g\equiv1$, use $C_k=\binom{n-1}{k-1}$, and otherwise use \eqref{eq:Ck-general} (DP with $A^{(0)}_{ij}\equiv 1$)\;
\BlankLine
\textbf{2. Run the sum-product dynamic program}\;
Compute $\widetilde{L}_{k j},\widetilde{R}_{k i}$ via \eqref{eq:LR} for $k=0,\dots,k_{\max}$ using $\widetilde{A}^{(0)}_{ij}$\;
Compute $P(k\mid y)\propto p(k)\,\widetilde{L}_{k n}/C_k$ via \eqref{eq:post-k}; set $\widehat{k}\leftarrow \arg\max_k P(k\mid y)$\;
\BlankLine
\textbf{3. Boundary marginals}\;
\For{$p=1$ \KwTo $\widehat{k}-1$}{
  For all $h\in\{p,\dots,n-(\widehat{k}-p)\}$, set
  $P(t_p=h\mid y,\widehat{k})\leftarrow \widetilde{L}_{p,h}\widetilde{R}_{\widehat{k}-p,h}/\widetilde{L}_{\widehat{k}n}$ (cf.\ \eqref{eq:boundary-post})\;
}
\BlankLine
\textbf{4. Export a joint MAP segmentation}\;
Run the max-sum/backpointer recursion of Section~\ref{sec:algorithms} using $\log \widetilde{A}^{(0)}_{ij}$ and the chosen count $\widehat{k}$\;
Backtrack to obtain $\widehat{t}^{\,\mathrm{MAP}}=(0=\widehat{t}^{\,\mathrm{MAP}}_0<\cdots<\widehat{t}^{\,\mathrm{MAP}}_{\widehat{k}}=n)$\;
\BlankLine
\textbf{5. Segment-level posterior moments on the exported segmentation}\;
\For{$q=1$ \KwTo $\widehat{k}$}{
  Let $(i,j]\leftarrow (\widehat{t}^{\,\mathrm{MAP}}_{q-1},\widehat{t}^{\,\mathrm{MAP}}_{q}]$\;
  For $r\in\{1,2\}$ set $\widehat{\mu}^{(r)}_q \leftarrow \widetilde{A}^{(r)}_{ij}/\widetilde{A}^{(0)}_{ij}$ (cf.\ \eqref{eq:segmom-seg})\;
}
\BlankLine
\textbf{6. Bayes regression curve within $\widehat{k}$}\;
Using $\widetilde{L},\widetilde{R},\widetilde{A}^{(r)}$ and \eqref{eq:segmom}, compute $\widehat{\mu}'_i$ and $\mathbb{V}[\mu'_i\mid y,\widehat{k}]$ for $i=1,\dots,n$\;
\BlankLine
\textbf{7. Optional: pooling/groups}\;
If multiple subjects or groups are present, pool block evidences (Theorem~\ref{thm:multisubject}) and/or run the latent-template EM routine of Section~\ref{sec:latent-em}, then repeat Steps 2--6.
\end{algorithm}


\subsection{Computational complexity and numerical stability}\label{sec:dp-complexity}
BayesBreak separates computation into a \emph{local} stage (building the block-evidence array) and a \emph{global} stage (running DP over partitions).
For conjugate exponential-family models, block evidences can be evaluated in $\mathcal{O}(1)$ amortized time per block once cumulative sufficient-statistic arrays are constructed, yielding $\mathcal{O}(n^2)$ total preprocessing.
The DP recursions over $(k,t)$ states are $\mathcal{O}(k_{\max} n^2)$ in the worst case for explicit segment-count posteriors; in many regimes, pruning and/or marginalization of $k$ can reduce the effective complexity, but BayesBreak's design prioritizes exactness and transparency of posterior quantities.
All forward/backward computations are performed in log-space with stable \texttt{logsumexp} operations, and implementations cache cumulative summaries (e.g., prefix sums of sufficient statistics) to avoid redundant work.
When storing the full forward table is memory-prohibitive, checkpoint layers or recompute forward messages on demand during backward passes, trading memory for additional passes through the DP.



\section{Irregular designs and length-aware priors}
\label{sec:irregular}

Irregularly spaced designs enter BayesBreak in two logically distinct places.
First, likelihood weights or exposures $w_t$ belong to the observation model when they encode
Poisson exposures, known Gaussian precisions, or Binomial trial counts. These quantities are
part of the sampling model; supply them only when substantively justified by the data-
generating process. Second, the geometry of the design points $x_1<\dots<x_n$ can enter through
a \emph{partition prior} that depends on physical segment lengths $\Delta_x(i,j)$. In most irregular-grid
applications, this second mechanism is the default one: irregular spacing by itself does not imply
that the likelihood is reweighted.

Accordingly, a clean default is to leave the observation model unchanged (often $w_t\equiv 1$) and
encode design geometry through the segment-factorized prior
\[
p(t\mid k) \propto \prod_{q=1}^k g(\Delta_x(t_{q-1},t_q)),
\]
with $g$ chosen to reflect prior beliefs about physical segment lengths. This is particularly natural
in genomics, spatial statistics, and event-time data, where the scientifically meaningful notion of
``distance'' is not the index increment $t_q-t_{q-1}$ but the physical length of the segment.

\begin{proposition}[Homogeneous Poisson-process prior: regular-grid index-uniform as a special case]
\label{prop:pp-index-uniform}
Partition the continuous interval $(u_0,u_n)$ into the $n$ candidate boundary intervals
$(u_{i-1},u_i]$ of lengths $\Delta u_i:=u_i-u_{i-1}$. Consider a homogeneous Poisson process
of rate $\lambda$ on $(u_0,u_n)$ and condition on the event that exactly $k-1$ events occur
and that no two events fall in the same inter-point interval. Then the induced discrete prior
on the selected set of $(k-1)$ boundary \emph{intervals} is proportional to the product of
their lengths, $\prod_{\ell\in\mathcal{I}} \Delta u_\ell$. In particular, on a regular grid
(all $\Delta u_\ell$ equal), the induced discrete prior is index-uniform.
\end{proposition}

\begin{proof}
For a homogeneous Poisson process, counts in disjoint intervals are independent and
$\#\{\text{events in }(u_{i-1},u_i]\}\sim\mathrm{Pois}(\lambda\Delta u_i)$. Conditioning on
the event that exactly one event occurs in each selected interval $\ell\in\mathcal{I}$ and
zero events occur in other intervals yields conditional probability proportional to
\[
\prod_{\ell\in\mathcal{I}} (\lambda\Delta u_\ell)e^{-\lambda\Delta u_\ell}
\prod_{\ell\notin\mathcal{I}} e^{-\lambda\Delta u_\ell}
\ \propto\ \prod_{\ell\in\mathcal{I}} \Delta u_\ell,
\]
since the global factor $e^{-\lambda\sum_\ell \Delta u_\ell}$ cancels. If $\Delta u_\ell$ are
all equal, all $(k-1)$-subsets $\mathcal{I}$ have equal probability, giving index-uniformity.
\end{proof}

\begin{remark}[Inhomogeneous intensity]
\label{rem:inhomog-pp}
No step of the argument uses that $\lambda$ is constant. For an absolutely continuous intensity
$\lambda(x)\ge 0$, the same conditioning yields a conditional prior proportional to
$\prod_{\ell\in\mathcal{I}}\Lambda_\ell$ where $\Lambda_\ell := \int_{x_{\ell-1}}^{x_\ell}\lambda(x)\,dx$
is the cumulative intensity on interval $\ell$. This gives a principled way to encode physically
motivated spatial variation in the prior hazard of a boundary.
\end{remark}

\begin{proposition}[Coarse-to-fine consistency of inherited partitions]
\label{thm:refine}
Let $x_m<\tilde x<x_{m+1}$ be a refinement of the design grid obtained by inserting a pseudo-index
between two existing design points but \emph{without} adding a new observation. Assume the pseudo-index
contributes no likelihood term and that block evidences are computed from the original observations only.
Then every coarse-grid segmentation $t$ has a lifted fine-grid segmentation $\tilde t$ that places no
boundary at the pseudo-index and satisfies
\[
\prod_{q} A^{(0)}_{t_{q-1}t_q}
  = \prod_q \widetilde{A}^{(0)}_{\tilde t_{q-1}\tilde t_q}
\]
for the likelihood contribution of the inherited partition. If, in addition, the partition prior is a
function of physical segment lengths $\Delta_x(t_{q-1},t_q)$, then the prior weight of this inherited
partition is unchanged as well.
\end{proposition}

\begin{proof}
Because no new observation is added, each coarse block $(t_{q-1},t_q]$ contains exactly the same data as
its lifted fine-grid counterpart $(\tilde t_{q-1},\tilde t_q]$. Their sufficient statistics and hence their
likelihood contributions coincide. If the prior depends only on physical segment lengths, then the lifted
block has the same endpoints in the physical domain and therefore the same prior factor.
\end{proof}

\begin{remark}[Why full refinement invariance is generally false]
The proposition above is deliberately weaker than a posterior invariance claim. Once the refined grid is
allowed to place \emph{new} boundaries at the inserted pseudo-index, the admissible partition space changes,
so the posterior distribution over segmentations and the normalization constants $C_k$ generally change as
well. In particular, refining the grid enlarges $\mathcal{T}_k$, which typically increases $C_k$ and
redistributes posterior mass over segmentations even when no new observations are added. Thus grid
refinement preserves the score of every \emph{inherited} partition, but it need not preserve the full
posterior unless one explicitly constrains the refined model to exclude new boundary locations.
\end{remark}

\begin{example}[Illustration on a $6$-point irregular design]
\label{ex:irregular-illustration}
Consider six observations at physical locations $x=(0,\,0.2,\,0.4,\,0.6,\,1.6,\,1.8)$ (a tight
cluster of four points followed by a pair after a large gap), with Gaussian responses having a
single genuine changepoint between $x=0.6$ and $x=1.6$. Under the index-uniform prior
$g\equiv 1$, the five candidate boundary indices are equally weighted a priori, so the posterior
must compete block-evidence differences against a flat prior. Under the length-aware prior
$g(\ell)\propto \ell$ (proportional to physical segment length, cf.\
Proposition~\ref{prop:pp-index-uniform}), the boundary index $i=4$ (separating the tight cluster
from the far pair) carries much larger prior mass, consistent with the physical intuition that
the prior favors a boundary at the large gap. This example is a qualitative prior-sensitivity illustration: changing $g$ shifts posterior
boundary mass without changing the likelihood. A larger quantitative sensitivity study is specified
as planned empirical work.
\end{example}

\section{Replicates and multi-subject pooling (shared boundaries)}
\label{sec:replicates}

\paragraph{Plate model (shared boundaries, subject-specific parameters).}
With a shared segmentation $t=(t_0,\dots,t_k)$ and subject-specific conjugate priors,
the generative model nests an inner \emph{segment} plate inside an outer \emph{subject} plate:
\begin{equation}
t\sim p(t\mid k),\qquad
\theta^{(s)}_q \stackrel{\mathrm{ind}}{\sim} p(\theta\mid \alpha^{(s)}_0,\beta^{(s)}_0),\quad
y^{(s)}_i\mid \theta^{(s)}_{q(i)},\,w^{(s)}_i \stackrel{\mathrm{ind}}{\sim} p(y\mid \theta^{(s)}_{q(i)},w^{(s)}_i),
\label{eq:plate-replicates}
\end{equation}
where $q(i)$ is the segment index containing observation $i$ under $t$.
\begin{figure}[H]
\centering
\begin{tikzpicture}
  \node[latent] (t) {$t$};
  \node[const, left=0.9cm of t] (pk) {$p(k)$};
  \node[latent, below=1.2cm of t] (th) {$\theta^{(s)}_q$};
  \node[obs, below=1.2cm of th] (y) {$y^{(s)}_i$};
  \node[const, left=0.9cm of th] (ab) {$\alpha^{(s)}_0,\beta^{(s)}_0$};
  \node[const, right=0.9cm of y] (w) {$w^{(s)}_i$};
  \edge {pk} {t};
  \edge {ab} {th};
  \edge {t,th,w} {y};
  \plate {pseg} {(th)} {$q=1{:}k$};
  \plate {psubj} {(pseg)(y)(ab)(w)} {$s=1{:}S$};
\end{tikzpicture}
\caption[Shared-boundary multi-subject plate model]{Plate model for multi-subject pooling under shared boundaries. The boundary vector $t$ is shared across subjects; each subject has its own segment parameters $\theta^{(s)}_q$ and exposures $w^{(s)}_i$.}
\label{fig:plate-replicates}
\end{figure}

Suppose we observe $S$ independent subjects, each with its own observations
$\{(x^{(s)}_i,y^{(s)}_i)\}$ and exposure weights $w^{(s)}_i$, but we wish to model them as
sharing a common segmentation (boundary locations) while allowing subject-specific segment
parameters. For simplicity of exposition, assume a common index grid $i=1,\dots,n$ (e.g., after
alignment/union of grids); subject-specific missing observations are handled by setting $w^{(s)}_i=0$
at missing indices.

Let $(i,j]$ denote a block and write subject-specific block sums as
$(S^{(s)}_{ij},W^{(s)}_{ij},H^{(s)}_{ij})$ with subject-specific prior hyperparameters
$(\alpha^{(s)}_0,\beta^{(s)}_0)$.

\begin{assumption}[Cross-subject conditional independence given shared boundaries]
\label{ass:cond-indep-subjects}
Conditional on a shared segmentation $t=(t_0,\dots,t_k)$, subjects $s=1,\dots,S$ are
independent, and each subject has its own segment parameters
$\theta^{(s)}_{1:k}$ with subject-specific conjugate priors.
\end{assumption}

\begin{theorem}[Exact pooling of block evidences]
\label{thm:multisubject}
Under Assumption~\ref{ass:cond-indep-subjects}, the pooled block evidence for $(i,j]$ factorizes as
\[
A^{(0)}_{ij,\mathrm{pool}}
 = \prod_{s=1}^S
    \exp\{H^{(s)}_{ij}\}\,
    \frac{Z(\alpha^{(s)}_0+S^{(s)}_{ij},\beta^{(s)}_0+W^{(s)}_{ij})}
         {Z(\alpha^{(s)}_0,\beta^{(s)}_0)}.
\]
Equivalently, $\log A^{(0)}_{ij,\mathrm{pool}}=\sum_s \log A^{(0)}_{ij,s}$. Running the DP
recursions \eqref{eq:LR}--\eqref{eq:segmom} with $A^{(0)}_{ij}$ replaced by
$A^{(0)}_{ij,\mathrm{pool}}$ yields the exact posterior over shared boundaries. Conditional on any
selected or averaged boundary configuration, the subject-specific segment-parameter posteriors are
then recovered from the corresponding subject-level conjugate updates.
\end{theorem}

\begin{proof}
\ProofStep{Step 1: subject-wise block evidences.}
For each subject $s$, Theorem~\ref{thm:ef-integral} yields the subject-specific block evidence
$A^{(0)}_{ij,s}$ computed from $(S^{(s)}_{ij},W^{(s)}_{ij},H^{(s)}_{ij})$.

\ProofStep{Step 2: independence implies product pooling.}
Conditional on shared boundaries, subjects are independent and have independent segment
parameters. Therefore, the joint likelihood (with segment parameters integrated out) for the
block is the product of the subject-wise block evidences:
$A^{(0)}_{ij,\mathrm{pool}}=\prod_s A^{(0)}_{ij,s}$, giving the stated form.

\ProofStep{Step 3: DP correctness under pooled evidences.}
The DP in Section~\ref{sec:dp} depends only on the block evidence matrix. Replacing
$A^{(0)}_{ij}$ by $A^{(0)}_{ij,\mathrm{pool}}$ therefore computes the exact shared-boundary
posterior for the pooled model; parameter posteriors remain available conditionally on boundaries
from the subject-specific block updates.
\end{proof}

\section{Groups with known membership}
\label{sec:groups-known}

\paragraph{Plate model (known groups, group-specific segmentations).}
Let $z_s\in\{1,\dots,G\}$ be the observed group label of subject $s$. The generative model nests
a group-specific boundary vector $t^{(g)}$ and group-specific segment parameters inside an outer
group plate, with subjects of each group drawn independently conditional on the corresponding
$t^{(g)}$:
\begin{equation}
t^{(g)}\sim p(t\mid k_g),\qquad
\theta^{(g,s)}_q \stackrel{\mathrm{ind}}{\sim} p(\theta\mid \alpha^{(g,s)}_0,\beta^{(g,s)}_0),\qquad
y^{(s)}_i\mid z_s=g, t^{(g)}, \theta^{(g,s)}_{q(i)} \stackrel{\mathrm{ind}}{\sim} p(y\mid \theta^{(g,s)}_{q(i)}, w^{(s)}_i).
\label{eq:plate-knowngroups}
\end{equation}
\begin{figure}[H]
\centering
\begin{tikzpicture}
  \node[obs] (z) {$z_s$};
  \node[latent, right=1.4cm of z] (tg) {$t^{(g)}$};
  \node[latent, below=1.1cm of tg] (thgs) {$\theta^{(g,s)}_q$};
  \node[obs, below=1.1cm of thgs] (y) {$y^{(s)}_i$};
  \node[const, right=0.9cm of y] (w) {$w^{(s)}_i$};
  \edge {tg} {thgs};
  \edge {thgs,w,z} {y};
  \plate {pg}   {(tg)} {$g=1{:}G$};
  \plate {pseg} {(thgs)} {$q=1{:}k_g$};
  \plate {psubj} {(pseg)(y)(z)} {$s=1{:}S$};
\end{tikzpicture}
\caption[Known-group plate model]{Plate model with known group memberships $z_s$. Each group has its own segmentation template $t^{(g)}$; subjects within a group share that template but retain subject-specific segment parameters $\theta^{(g,s)}_q$.}
\label{fig:plate-knowngroups}
\end{figure}

In some applications subjects can be partitioned into groups $g=1,\dots,G$ (e.g.\ experimental
conditions), and we allow each group to have its own segmentation while still
pooling within groups. When group memberships are known, BayesBreak handles this case by
applying the pooled DP of Section~\ref{sec:replicates} independently within each group.

Concretely, for group $g$ with subject index set $\mathcal{S}_g$, we pool per-subject block
evidences $\{A^{(0)}_{ij,s}\}_{s\in\mathcal{S}_g}$ into a group-specific pooled evidence
$A^{(0)}_{ij,g}$ using Theorem~\ref{thm:multisubject}, then run the DP recursions
\eqref{eq:LR}--\eqref{eq:segmom} with $A^{(0)}_{ij}=A^{(0)}_{ij,g}$. This yields group-specific
$P_g(k\mid y)$, boundary posteriors $P_g(t_p\mid y,k)$, and segment/regression-curve moments.
The total computational cost is $G$ times the single-group cost, since DP runs are independent
across groups conditional on the observed labels; block-evidence matrices can share a common
prefix-sum precomputation only when the underlying design, weights, and missingness pattern permit
it. This known-group construction is the supervised counterpart of the latent-template mixture in
Section~\ref{sec:latent-em}.

\section{Latent groups and an EM algorithm}
\label{sec:latent-em}

\paragraph{Plate model (latent-group template mixture).}
Group labels $z_s\in\{1,\dots,G\}$ are now \emph{latent}, with Categorical prior $\pi$. Each
group $g$ carries its own template $\tau_g=(k_g,t^{(g)})$; subject-specific segment parameters
$\theta^{(g,s)}_q$ are integrated out inside block evidences. The generative model is
\begin{equation}
\begin{aligned}
\pi&\sim p(\pi),\quad \tau_g\sim p(\tau),\quad z_s\mid\pi\stackrel{\mathrm{ind}}{\sim}\mathrm{Cat}(\pi),\\
\theta^{(g,s)}_q&\stackrel{\mathrm{ind}}{\sim}p(\theta\mid\cdot),\qquad
y^{(s)}_i\mid z_s,\tau_{z_s},\theta^{(z_s,s)}_{q(i)}\stackrel{\mathrm{ind}}{\sim}p(y\mid\theta^{(z_s,s)}_{q(i)},w^{(s)}_i).
\end{aligned}
\label{eq:plate-latent-em}
\end{equation}
\begin{figure}[H]
\centering
\begin{tikzpicture}
  \node[latent] (pi) {$\pi$};
  \node[latent, right=1.4cm of pi] (tau) {$\tau_g$};
  \node[latent, below=1.1cm of pi] (z) {$z_s$};
  \node[latent, below=1.1cm of tau] (thgs) {$\theta^{(g,s)}_q$};
  \node[obs, below=1.1cm of thgs] (y) {$y^{(s)}_i$};
  \node[const, right=0.9cm of y] (w) {$w^{(s)}_i$};
  \edge {pi} {z};
  \edge {tau} {thgs};
  \edge {z,tau,thgs,w} {y};
  \plate {pg}   {(tau)} {$g=1{:}G$};
  \plate {pseg} {(thgs)} {$q=1{:}k_g$};
  \plate {psubj} {(pseg)(y)(z)} {$s=1{:}S$};
\end{tikzpicture}
\caption[Latent-group template-mixture plate model]{Plate model for the latent-group template mixture. Unshaded $z_s$ are latent; EM alternates between exact E-step responsibilities $p(z_s=g\mid y^{(s)},\pi,\tau)$ and M-step max-sum updates of the templates $\tau_g$.}
\label{fig:plate-latent-em}
\end{figure}

For unknown group membership we use a \emph{latent template-mixture} model. Each latent group carries a
single segmentation template that is shared across the subjects assigned to that group, while subject-
specific segment parameters are still integrated out inside each block evidence. This is the cleanest
setting in which EM updates remain exact and interpretable.

Let $\tau_g=(k_g,t^{(g)})$ denote the template of group $g$, where
$1\le k_g\le k_{\max}$ and $0=t^{(g)}_0<\cdots<t^{(g)}_{k_g}=n$. Conditional on membership in group $g$,
subject $s$ is assigned the positive prior-weighted template score
\begin{equation}
S_g\big(y^{(s)};\tau_g\big)
 = \frac{p(k_g)}{C_{k_g}}
   \prod_{q=1}^{k_g} \widetilde{A}^{(0,s)}_{t^{(g)}_{q-1}t^{(g)}_q},
\label{eq:template-lik}
\end{equation}
where $\widetilde{A}^{(0,s)}_{ij}=A^{(0,s)}_{ij}g(\Delta_x(i,j))$ is the subject-specific block evidence with any
length factor already absorbed. This score defines the finite template-mixture objective optimized
below; it is distinct from a fully Bayesian posterior over all group segmentations. The
group mixture weights are nonnegative with $\sum_g \pi_g=1$; implementations either deactivate
empty groups with $\pi_g=0$ or use a small Dirichlet pseudocount/floor to keep every group active.

Introduce latent indicators $z_{sg}\in\{0,1\}$ with $\sum_g z_{sg}=1$. The complete-data log-likelihood is
\[
\mathcal{J}(z,\pi,\tau)
 = \sum_{s=1}^S\sum_{g=1}^G z_{sg}\bigl[\log\pi_g + \log S_g(y^{(s)};\tau_g)\bigr].
\]
Let $r_{sg}:=\mathbb{E}[z_{sg}\mid y,\pi,\tau]$ denote the usual responsibilities.

\paragraph{E-step.}
Given current $(\pi,\tau)$,
\begin{equation}
r_{sg}
 = \frac{\pi_g\,S_g(y^{(s)};\tau_g)}{\sum_{g'=1}^G \pi_{g'}\,S_{g'}(y^{(s)};\tau_{g'})}.
\label{eq:template-resp}
\end{equation}

\paragraph{M-step for mixing weights.}
For fixed responsibilities,
\[
\pi_g \leftarrow \frac{1}{S}\sum_{s=1}^S r_{sg}.
\]

\paragraph{M-step for group templates.}
For fixed $r$, write $n_g(r):=\sum_{s=1}^S r_{sg}$. The $\tau_g$-dependent part of the EM
auxiliary function is
\[
Q_g(\tau_g)
 = \sum_{s=1}^S r_{sg}\log S_g(y^{(s)};\tau_g)
 = n_g(r)\{\log p(k_g)-\log C_{k_g}\}
   + \sum_{q=1}^{k_g} B^{(g)}_{t^{(g)}_{q-1}t^{(g)}_q},
\]
where the responsibility-weighted block score is
\begin{equation}
B^{(g)}_{ij}
 := n_g(r)\log g(\Delta_x(i,j)) + \sum_{s=1}^S r_{sg}\,\log A^{(0,s)}_{ij}.
\label{eq:template-blockscore}
\end{equation}
Thus the template-prior normalizer and the length factor are scaled by the same effective group
size $n_g(r)$ as the rest of the prior-weighted score. For each group $g$, maximizing
$Q_g(\tau_g)$ is then exactly the same max-sum segmentation problem as in
Section~\ref{sec:algorithms}, except that the block score is replaced by $B^{(g)}_{ij}$ and the
count-specific offset $n_g(r)\{\log p(k)-\log C_k\}$ is added when selecting $k_g$. In other
words, the M-step is an exact responsibility-weighted backtracking DP for the stated finite
objective.

\begin{theorem}[Monotone EM updates for the template mixture]
\label{thm:em-monotone}
For the finite template-mixture objective defined by \eqref{eq:template-lik}, the EM iteration
consisting of, in order,
(i) the exact responsibility update \eqref{eq:template-resp} given current $(\pi^{(t)},\tau^{(t)})$
to produce $r^{(t+1)}$;
(ii) the closed-form mixing-weight update
$\pi^{(t+1)} = \arg\max_\pi Q(\pi,\tau^{(t)};r^{(t+1)})$; and
(iii) the exact max-sum DP update of each template
$\tau^{(t+1)} = \arg\max_\tau Q(\pi^{(t+1)},\tau;r^{(t+1)})$ via \eqref{eq:template-blockscore}
does not decrease the finite-mixture objective
\[
\ell_\star(\pi,\tau) = \sum_{s=1}^S \log\Big(\sum_{g=1}^G \pi_g S_g(y^{(s)};\tau_g)\Big).
\]
Any fixed point of these updates, with deterministic tie-breaking for equal-scoring templates, is a
coordinatewise maximizer of the EM auxiliary function. Across restarts, the selected solution is the
run with the largest achieved finite-mixture objective.
\end{theorem}

\begin{proof}
The E-step with current $(\pi^{(t)},\tau^{(t)})$ computes the exact conditional responsibilities
for the current finite set of template scores, which maximizes the standard EM lower bound for
$\ell_\star$ with respect to the assignment distribution $q$ for fixed parameters. The M-step
maximizes the auxiliary function
$Q(\pi,\tau;r^{(t+1)}):=\sum_{s,g}r^{(t+1)}_{sg}[\log\pi_g+\log S_g(y^{(s)};\tau_g)]$ in two
coordinates: (ii) over $\pi\in\Delta^{G-1}$ (closed-form simplex maximizer
$\pi_g^{(t+1)}=S^{-1}\sum_s r^{(t+1)}_{sg}$ when no pseudocount is used), and (iii) over each
$\tau_g$ independently. The latter is an additive max-sum segmentation objective with block score
\eqref{eq:template-blockscore} and count offset $n_g(r)\{\log p(k)-\log C_k\}$, exactly
maximizable by DP and backtracking. Since each coordinate update exactly maximizes its factor,
$Q(\pi^{(t+1)},\tau^{(t+1)};r^{(t+1)})\ge Q(\pi^{(t)},\tau^{(t)};r^{(t+1)})$, and the finite-mixture
objective $\ell_\star$ is non-decreasing by the generalized EM monotonicity theorem
\citep{dempster1977em}. A fixed point satisfies the coordinatewise maximization conditions for the
auxiliary function, rather than differentiable stationarity in a continuous template space.
\end{proof}

\paragraph{Implementation.}
Algorithm~\ref{alg:multi-em} emphasizes the separation between
(i) exact responsibility updates that use current group templates and
(ii) exact template updates obtained from responsibility-weighted block scores.
Because the template space is discrete while the mixture weights are continuous, multiple local optima occur. Use at least $R\ge 5$ restarts from either (a) a $k$-means clustering of the
per-subject pooled block-evidence matrices, (b) templates drawn from the partition prior, or
(c) a warm start derived from a known-groups solution on a small labeled subset when available.
Declare convergence when (i) no template changes between two successive iterations,
(ii) the finite-mixture objective change falls below a user tolerance (e.g., $10^{-6}$ relative),
and (iii) deterministic tie-breaking would choose the same templates if the weighted score matrix
were recomputed. Template stability alone is a useful diagnostic but does not by itself certify
that all objective and tie-handling conditions have been met.

\begin{algorithm}[H]
\caption{Latent-template EM for grouped segmentation (log-space)}
\label{alg:multi-em}
\KwIn{Per-subject log block evidences $\{\ell^{(0,s)}_{ij}:=\log A^{(0,s)}_{ij}\}_{s=1}^S$; number of groups $G$; $k_{\max}$; prior $\log p(k)$; length factor $g$; tolerance $\epsilon$}
Initialize mixture weights $\{\pi_g\}_{g=1}^G$ and template segmentations $\{\tau_g=(k_g,t^{(g)})\}_{g=1}^G$ (several restarts; cf.\ initialization guidance above)\;
\Repeat{template vector stable AND $|\ell_\star^{(t+1)}-\ell_\star^{(t)}|<\epsilon(1+|\ell_\star^{(t)}|)$}{
  \textbf{E-step (log-space; order: responsibilities first):}\;
  \For{$s=1$ \KwTo $S$}{
    \For{$g=1$ \KwTo $G$}{
      $\log S_g(y^{(s)};\tau_g) \leftarrow \log p(k_g) - \log C_{k_g} + \sum_{q=1}^{k_g}\ell^{(0,s)}_{t^{(g)}_{q-1}t^{(g)}_q} + \sum_{q=1}^{k_g}\log g(\Delta_x(t^{(g)}_{q-1},t^{(g)}_q))$\;
      $\log u_{sg}\leftarrow \log\pi_g + \log S_g(y^{(s)};\tau_g)$\;
    }
    $\log r_{sg}\leftarrow \log u_{sg} - \texttt{logsumexp}_{g'}(\log u_{sg'})$;\quad $r_{sg}\leftarrow \exp(\log r_{sg})$\;
  }
  Accumulate finite-mixture objective $\ell_\star^{(t+1)}\leftarrow \sum_s \texttt{logsumexp}_g(\log u_{sg})$\;
  \textbf{M-step (in the order weights then templates):}\;
  $\pi_g\leftarrow S^{-1}\sum_s r_{sg}$ for each $g$\;
  \For{$g=1$ \KwTo $G$}{
    Set $n_g\leftarrow\sum_s r_{sg}$ and form scores $B^{(g)}_{ij} = n_g\log g(\Delta_x(i,j)) + \sum_{s=1}^S r_{sg}\,\ell^{(0,s)}_{ij}$\;
    Run max-sum DP/backtracking over $k\in\{1,\dots,k_{\max}\}$ using score matrix $B^{(g)}$ and offset $n_g\{\log p(k)-\log C_k\}$\;
    Update $\tau_g=(k_g,t^{(g)})$ to the maximizing template\;
  }
}
\KwOut{$\{\pi_g\}$, $\{\tau_g\}$, and responsibilities $\{r_{sg}\}$}
\end{algorithm}

\section{Family-specific derivations}
\label{sec:families}

We now instantiate \eqref{eq:A0}--\eqref{eq:Ar} for three closed-form conjugate families
(Gaussian, Poisson, and Binomial), an overdispersed count model with a conjugate distributional
parameter, and one low-dimensional quadrature example for $(0,1)$-valued responses. In each case
we give (i) the segment model and prior, (ii) posterior hyperparameters, (iii) the block evidence
$A^{(0)}_{ij}$, and (iv) first two posterior moments on the scale used by the Bayes-curve
interface. For non-Gaussian families this scale must be stated explicitly: natural-parameter,
distribution-parameter, and observation-mean moments are not interchangeable.
Implementation routines are provided immediately after each derivation
(Algorithms~\ref{alg:gaussian-block}--\ref{alg:betaobs-block}).

\subsection{Gaussian with known variance (heteroscedastic via weights)}
Assume
\[
y_t\mid \mu \sim \mathcal{N}\!\big(\mu,\sigma^2/w_t\big),\qquad
\mu\sim \mathcal{N}(\nu,\rho^2),
\]
with known $\sigma^2$ and positive weights $w_t$ (e.g.\ inverse known variances or exposure).
For a block $(i,j]$, define
\[
\begin{aligned}
n_{ij}&:=j-i, &
W_{ij}&:=\sum_{t=i+1}^j w_t, &
S_{ij}&:=\sum_{t=i+1}^j w_t y_t,\\
Q_{ij}&:=\sum_{t=i+1}^j w_t (y_t-\nu)^2, &
\kappa_{ij}&:=\rho^2 W_{ij}+\sigma^2.
\end{aligned}
\]
Then
\[
\mu\mid(i,j] \sim \mathcal{N}\!\Big(
\frac{\rho^2 S_{ij} + \sigma^2 \nu}{\kappa_{ij}},
\ \frac{\sigma^2\rho^2}{\kappa_{ij}}
\Big).
\]
The integrated block evidence is
\begin{align*}
A^{(0)}_{ij}
&= \int \Big[\prod_{t=i+1}^j \mathcal{N}(y_t\mid \mu,\sigma^2/w_t)\Big]
\mathcal{N}(\mu\mid \nu,\rho^2)\,d\mu\\
&=(2\pi\sigma^2)^{-n_{ij}/2}
\Big(\prod_{t=i+1}^j w_t^{1/2}\Big)
\Big(1+\tfrac{\rho^2}{\sigma^2}W_{ij}\Big)^{-1/2}
\exp\!\left\{-\frac{Q_{ij}}{2\sigma^2}
+\frac{\big(S_{ij}-\nu W_{ij}\big)^2}
{2\big(\sigma^2 W_{ij}+\sigma^4/\rho^2\big)}\right\}.
\end{align*}
Consequently,
\[
A^{(1)}_{ij} = A^{(0)}_{ij}\,\mathbb{E}[\mu\mid(i,j]],
\qquad
A^{(2)}_{ij} = A^{(0)}_{ij}\,\Big(\mathbb{V}[\mu\mid(i,j]] + \mathbb{E}[\mu\mid(i,j]]^2\Big).
\]


\paragraph{Implementation.}
Algorithm~\ref{alg:gaussian-block} gives a numerically stable implementation of the Gaussian block log-evidence and posterior moments using cumulative weighted sums.
The routine is written so that the DP layer consumes only $\log \mathcal{L}_{ij}$ and does not depend on Gaussian-specific algebra.

\begin{algorithm}[H]
\caption{GaussianBlock$(y,w,i,j,\nu,\rho^2,\sigma^2)$}
\label{alg:gaussian-block}
\KwIn{Block indices $(i,j]$; observations $y_{i+1:j}$; weights $w_{i+1:j}$}
$n\leftarrow j-i$;\quad $W\leftarrow \sum_{t=i+1}^j w_t$;\quad $S\leftarrow \sum_{t=i+1}^j w_t y_t$;\quad $Q\leftarrow \sum_{t=i+1}^j w_t (y_t-\nu)^2$\;
$\kappa \leftarrow \rho^2 W+\sigma^2$\;
$\log A0 \leftarrow -\tfrac{1}{2}n\log(2\pi)
-\tfrac{1}{2}\sum_{t=i+1}^j \log(\sigma^2/w_t)
-\tfrac{Q}{2\sigma^2}
+\tfrac{(S-\nu W)^2}{2(\sigma^2 W+\sigma^4/\rho^2)}
-\tfrac{1}{2}\log\big(1+\tfrac{\rho^2}{\sigma^2}W\big)$\;
$A0\leftarrow \exp(\log A0)$\;
$\mu_{\text{mean}}\leftarrow (\rho^2 S+\sigma^2\nu)/\kappa$;\quad
$\mu_{\text{var}}\leftarrow \sigma^2\rho^2/\kappa$\;
\KwOut{$A0$, $\mu_{\text{mean}}$, $\mu_{\text{var}}$}
\end{algorithm}


\subsection{Poisson with exposure}
Assume a blockwise constant rate $\lambda$ with exposure weights $w_t>0$:
\[
y_t\mid\lambda \sim \mathrm{Pois}(\lambda w_t),\qquad
\lambda\sim\mathrm{Gamma}(a_0,b_0)\quad\text{(shape--rate)}.
\]
For a block $(i,j]$, define
\[
C_{ij} := \sum_{t=i+1}^j y_t,\qquad
W_{ij} := \sum_{t=i+1}^j w_t,\qquad
H_{ij} := \sum_{t=i+1}^j \Big[y_t\log w_t - \log(y_t!)\Big].
\]
Then
\[
\lambda\mid(i,j] \sim \mathrm{Gamma}(a_0+C_{ij},\,b_0+W_{ij}),
\]
so
\[
\mathbb{E}[\lambda\mid(i,j]] = \frac{a_0+C_{ij}}{b_0+W_{ij}},\qquad
\mathbb{V}[\lambda\mid(i,j]] = \frac{a_0+C_{ij}}{(b_0+W_{ij})^2}.
\]
The integrated block evidence is
\[
A^{(0)}_{ij}
= \exp\{H_{ij}\}\,
\frac{b_0^{a_0}\,\Gamma(a_0+C_{ij})}{\Gamma(a_0)\,(b_0+W_{ij})^{a_0+C_{ij}}}.
\]
Moments satisfy the usual relation
\[
A^{(1)}_{ij} = A^{(0)}_{ij}\,\mathbb{E}[\lambda\mid(i,j]],\qquad
A^{(2)}_{ij} = A^{(0)}_{ij}\,\Big(\mathbb{V}[\lambda\mid(i,j]] + \mathbb{E}[\lambda\mid(i,j]]^2\Big).
\]


\paragraph{Implementation.}
Algorithm~\ref{alg:poisson-block} summarizes the Poisson-with-exposure block evidence computation (and associated posterior summaries) in a form that matches the BayesBreak block-evidence interface.

\begin{algorithm}[H]
\caption{PoissonBlock$(y,w,i,j,a_0,b_0)$}
\label{alg:poisson-block}
\KwIn{Block $(i,j]$; counts $y_{i+1:j}$; exposures $w_{i+1:j}$}
$C\leftarrow \sum_{t=i+1}^j y_t$;\quad $W\leftarrow \sum_{t=i+1}^j w_t$\;
$\log A0 \leftarrow \sum_{t=i+1}^j y_t\log w_t
- \sum_{t=i+1}^j \log(y_t!)
+ a_0\log b_0
+ \log\Gamma(a_0+C)
- \log\Gamma(a_0)
- (a_0+C)\log(b_0+W)$\;
$A0\leftarrow \exp(\log A0)$\;
$\lambda_{\text{mean}}\leftarrow (a_0+C)/(b_0+W)$;\quad
$\lambda_{\text{var}}\leftarrow (a_0+C)/(b_0+W)^2$\;
\KwOut{$A0$, $\lambda_{\text{mean}}$, $\lambda_{\text{var}}$}
\end{algorithm}


\subsection{Binomial with Beta prior}
Assume blockwise constant success probability $p$ with heterogeneous trial counts $m_t$:
\[
y_t\mid p \sim\mathrm{Binom}(m_t,p),\qquad
p\sim\mathrm{Beta}(a_0,b_0).
\]
For a block $(i,j]$, define
\[
C_{ij} := \sum_{t=i+1}^j y_t,\qquad
M_{ij} := \sum_{t=i+1}^j m_t,\qquad
H_{ij} := \sum_{t=i+1}^j \log\binom{m_t}{y_t}.
\]
Then $p\mid(i,j]\sim\mathrm{Beta}(a_0+C_{ij}, b_0+M_{ij}-C_{ij})$, so
\[
\mathbb{E}[p\mid(i,j]]=\frac{a_0+C_{ij}}{a_0+b_0+M_{ij}},\quad
\mathbb{V}[p\mid(i,j]]=
\frac{(a_0+C_{ij})(b_0+M_{ij}-C_{ij})}
{(a_0+b_0+M_{ij})^2(a_0+b_0+M_{ij}+1)}.
\]
The block evidence is
\[
A^{(0)}_{ij}
= \exp\{H_{ij}\}
\frac{B(a_0+C_{ij},\,b_0+M_{ij}-C_{ij})}{B(a_0,b_0)}.
\]
Again,
\[
A^{(1)}_{ij} = A^{(0)}_{ij}\,\mathbb{E}[p\mid(i,j]],\qquad
A^{(2)}_{ij} = A^{(0)}_{ij}\,\Big(\mathbb{V}[p\mid(i,j]] + \mathbb{E}[p\mid(i,j]]^2\Big).
\]


\paragraph{Implementation.}
Algorithm~\ref{alg:binomial-block} provides pseudocode for Binomial/Bernoulli block evidence evaluation under Beta conjugacy, along with posterior mean/variance computations used for segment-wise moments and MAP/Bayes curve summaries.

\begin{algorithm}[H]
\caption{BinomialBlock$(y,m,i,j,a_0,b_0)$}
\label{alg:binomial-block}
\KwIn{Block $(i,j]$; successes $y_{i+1:j}$; trials $m_{i+1:j}$}
$C\leftarrow \sum_{t=i+1}^j y_t$;\quad $M\leftarrow \sum_{t=i+1}^j m_t$\;
$\log A0 \leftarrow \sum_{t=i+1}^j \log \binom{m_t}{y_t}
+ \log B(a_0+C,b_0+M-C)
- \log B(a_0,b_0)$\;
$A0\leftarrow \exp(\log A0)$\;
$p_{\text{mean}}\leftarrow (a_0+C)/(a_0+b_0+M)$\;
$p_{\text{var}}\leftarrow \frac{(a_0+C)(b_0+M-C)}
{(a_0+b_0+M)^2(a_0+b_0+M+1)}$\;
\KwOut{$A0$, $p_{\text{mean}}$, $p_{\text{var}}$}
\end{algorithm}

\subsection{Negative-Binomial block (overdispersed counts) with fixed dispersion}
\label{sec:nb-block}
For count data with extra-Poisson variability, a natural block model fixes a dispersion parameter
$r>0$ and treats the success probability $p\in(0,1)$ per segment with a conjugate Beta prior:
\[
y_t\mid p \sim \mathrm{NegBin}(r, p),\qquad p\sim\mathrm{Beta}(a_0,b_0),
\]
where $\mathrm{NegBin}(r,p)$ has pmf $\binom{y+r-1}{y}(1-p)^y p^r$ and count mean
$r(1-p)/p$. For a block $(i,j]$, set $C_{ij}:=\sum_{t=i+1}^j y_t$, $N_{ij}:=\sum_{t=i+1}^j r_t$
(with $r_t\equiv r$ in the fixed-dispersion case), and
$H_{ij}:=\sum_{t=i+1}^j \log\binom{y_t+r_t-1}{y_t}$. Beta--NegBin conjugacy gives
$p\mid(i,j]\sim\mathrm{Beta}(a_B,b_B)$ with $a_B=a_0+N_{ij}$ and $b_B=b_0+C_{ij}$, so
\[
A^{(0)}_{ij}=\exp\{H_{ij}\}\,\frac{B(a_B,b_B)}{B(a_0,b_0)}.
\]
The Beta posterior moments of $p$ are useful diagnostics, but the Bayes-curve target on the
observation scale is a count mean. For a future count with dispersion $r_\ast$, this target is
$m_\ast(p)=r_\ast(1-p)/p$, giving
\begin{align*}
\mathbb{E}[m_\ast(p)\mid(i,j]]
&=r_\ast\frac{b_B}{a_B-1}\quad(a_B>1),\\
\mathbb{V}[m_\ast(p)\mid(i,j]]
&=r_\ast^2\frac{b_B(a_B+b_B-1)}{(a_B-1)^2(a_B-2)}\quad(a_B>2).
\end{align*}
Moment numerators $A^{(1)}_{ij}$ and $A^{(2)}_{ij}$ therefore multiply $A^{(0)}_{ij}$ by
the corresponding observation-mean moments, not by raw Beta moments of $p$, unless the exported
quantity is explicitly labeled as parameter-scale. When $r$ is unknown, a semi-conjugate EM or a
one-dimensional quadrature in $r$ can be layered on top; we omit the details.

\paragraph{Implementation.}
Algorithm~\ref{alg:nb-block} summarizes the Negative-Binomial block computation in the same
interface as Algorithms~\ref{alg:gaussian-block}--\ref{alg:binomial-block}.

\begin{algorithm}[H]
\caption{NegBinBlock$(y,r,i,j,a_0,b_0)$}
\label{alg:nb-block}
\KwIn{Block $(i,j]$; counts $y_{i+1:j}$; dispersion $r$ (possibly per-$t$)}
$C\leftarrow\sum_{t=i+1}^j y_t$;\quad $N\leftarrow\sum_{t=i+1}^j r_t$\;
$\log A0\leftarrow \sum_{t=i+1}^j \log\binom{y_t+r_t-1}{y_t} + \log B(a_0+N, b_0+C) - \log B(a_0,b_0)$\;
$A0\leftarrow\exp(\log A0)$\;
$a_B\leftarrow a_0+N$;\quad $b_B\leftarrow b_0+C$\;
$\texttt{mean\_per\_disp}\leftarrow b_B/(a_B-1)$ if $a_B>1$;\quad
$\texttt{var\_per\_disp}\leftarrow b_B(a_B+b_B-1)/\{(a_B-1)^2(a_B-2)\}$ if $a_B>2$\;
\KwOut{$A0$, posterior summaries for $p$, and observation-mean moments $r_\ast\times\texttt{mean\_per\_disp}$ when a prediction dispersion $r_\ast$ is specified}
\end{algorithm}


\subsection{Beta-distributed observations with known precision}
Assume $y_t\in(0,1)$ with latent mean $\mu$ constant within a block and known precision
$\phi_t>0$ (with $\phi_t\equiv\phi$ as the fixed-precision special case):
\[
y_t\mid \mu \sim \mathrm{Beta}(\phi_t\mu,\phi_t(1-\mu)),\qquad
\mu\sim\mathrm{Beta}(a_0,b_0).
\]
Conjugacy in $\mu$ is not available, but all integrals remain one-dimensional and can be computed
by deterministic quadrature. We do not require or claim global log-concavity of the induced
posterior in $\mu$; instead, implementations check quadrature stability by increasing the
node count and by controlling behavior near the endpoints of $(0,1)$. For a block $(i,j]$ define
\[
\ell_{ij}(\mu)
:= \sum_{t=i+1}^j \log \mathrm{BetaPDF}\big(y_t\mid \phi_t\mu,\phi_t(1-\mu)\big),
\qquad
\log\pi(\mu) := \log \mathrm{BetaPDF}(\mu\mid a_0,b_0).
\]
Then for $r=0,1,2$,
\[
A^{(r)}_{ij} = \int_0^1 \mu^{\,r}\,\exp\big\{\ell_{ij}(\mu)+\log\pi(\mu)\big\}\,d\mu.
\]
In practice, approximate via Gauss--Legendre (or similar) quadrature with nodes $\{\mu_g\}$ and
weights $\{w_g\}$:
\[
A^{(r)}_{ij} \approx \sum_{g=1}^G w_g\,\mu_g^{\,r}\,
\exp\big\{\ell_{ij}(\mu_g)+\log\pi(\mu_g)\big\}.
\]


\paragraph{Implementation.}
Algorithm~\ref{alg:betaobs-block} implements this one-dimensional quadrature routine. It is useful for continuous proportions and other $(0,1)$-valued signals, and is a deterministic non-conjugate block integral rather than a conjugate closed form.

\begin{algorithm}[H]
\caption{BetaObsBlock$(y,\phi,i,j,a_0,b_0,\{\mu_g,w_g\}_{g=1}^G)$}
\label{alg:betaobs-block}
\KwIn{Block $(i,j]$; observations $y_{i+1:j}\in(0,1)$; known precisions $\phi_{i+1:j}$; Beta prior $(a_0,b_0)$; quadrature nodes $\{\mu_g,w_g\}$}
$S_0\leftarrow 0$, $S_1\leftarrow 0$, $S_2\leftarrow 0$\;
\For{$g=1$ \KwTo $G$}{
  $\log \ell_g \leftarrow \sum_{t=i+1}^j \log \mathrm{BetaPDF}\big(y_t\mid \phi_t\mu_g,\phi_t(1-\mu_g)\big)$\;
  $\log \pi_g \leftarrow \log \mathrm{BetaPDF}(\mu_g\mid a_0,b_0)$\;
  $v_g \leftarrow \exp(\log \ell_g + \log \pi_g)$\;
  $S_0 \leftarrow S_0 + w_g v_g$;\quad
  $S_1 \leftarrow S_1 + w_g v_g \mu_g$;\quad
  $S_2 \leftarrow S_2 + w_g v_g \mu_g^2$\;
}
$A0\leftarrow S_0$;\quad
$\mu_{\text{mean}}\leftarrow S_1/S_0$;\quad
$\mu_{\text{var}}\leftarrow S_2/S_0 - (\mu_{\text{mean}})^2$\;
\KwOut{$A0$, $\mu_{\text{mean}}$, $\mu_{\text{var}}$}
\end{algorithm}

\subsection{Summary of block families}
\label{sec:families-summary}
Table~\ref{tab:family-summary} collects the block families derived in this section in a single
reference view. The conjugate rows expose closed-form ratios of normalizers, while the
Beta-observation row supplies deterministic quadrature values. In both cases the DP layer of
Section~\ref{sec:dp} only requires a compatible block-evidence array and is invariant to the
family used to fill that array.

\begin{table}[H]
\centering
\small
\resizebox{\textwidth}{!}{%
\begin{tabular}{@{}lllll@{}}
\toprule
Family & Prior & Posterior & Block evidence $A^{(0)}_{ij}$ & Mean $\mathbb{E}[m(\theta)|(i,j]]$ \\
\midrule
Gaussian$(\mu,\sigma^2/w)$ & $\mathcal{N}(\nu,\rho^2)$ & $\mathcal{N}(\nu_B,\rho_B^2)$ & closed-form normal ratio & $(\rho^2 S+\sigma^2\nu)/\kappa$ \\
Poisson$(\lambda w)$ & $\mathrm{Gamma}(a_0,b_0)$ & $\mathrm{Gamma}(a_0+C,b_0+W)$ & $\propto b_0^{a_0}\Gamma(a_0+C)/(b_0+W)^{a_0+C}$ & $(a_0+C)/(b_0+W)$ \\
Binomial$(m,p)$ & $\mathrm{Beta}(a_0,b_0)$ & $\mathrm{Beta}(a_0+C,b_0+M-C)$ & $B(a_0+C,b_0+M-C)/B(a_0,b_0)$ & $(a_0+C)/(a_0+b_0+M)$ \\
NegBin$(r,p)$ & $\mathrm{Beta}(a_0,b_0)$ & $\mathrm{Beta}(a_B,b_B)$ & $B(a_B,b_B)/B(a_0,b_0)$ & $r_\ast b_B/(a_B-1)$ \\
Beta-obs$(\mu,\phi_t)$ & $\mathrm{Beta}(a_0,b_0)$ & (1D quadrature) & Gauss--Legendre nodes & from quadrature \\
\bottomrule
\end{tabular}%
}
\caption{Block families derived in Section~\ref{sec:families}. Base-measure factor
$\exp\{H_{ij}\}$ is implicit in the conjugate rows. The first three rows are closed-form
conjugate examples; the Negative-Binomial row distinguishes parameter-scale Beta moments from the
observation-mean target; and the Beta-observation row is non-conjugate but remains compatible with
the DP layer through one-dimensional quadrature.}
\label{tab:family-summary}
\end{table}

\subsection{Posterior-predictive formulas for common conjugate families (optional)}
\label{sec:families-predictive}
For groups that export segment-level conjugate posteriors, the segment posterior-predictive
\eqref{eq:pp-seg} reduces to standard closed forms:

\paragraph{Gaussian (known variance).}
If a segment posterior is $\mu\mid \text{train}\sim \mathcal{N}(\nu_B,\rho_B^2)$ and
$y\mid\mu\sim\mathcal{N}(\mu,\sigma^2/w)$, then the posterior-predictive for a new point is
Gaussian: $y\mid\text{train}\sim \mathcal{N}(\nu_B,\rho_B^2+\sigma^2/w)$. For $M$ new
points in the segment with weights $w_1,\dots,w_M$, the joint posterior-predictive is the
multivariate Gaussian
\[
\mathbf{y}^{\mathrm{new}}\mid\text{train}\sim \mathcal{N}_M\big(\nu_B\mathbf{1}_M,\ \rho_B^2\,\mathbf{1}_M\mathbf{1}_M^\top + \mathrm{diag}(\sigma^2/w_1,\dots,\sigma^2/w_M)\big),
\]
which has rank-one-plus-diagonal covariance and whose log density can be evaluated in
$\mathcal{O}(M)$ by the Sherman--Morrison identity.

\paragraph{Poisson--Gamma.}
If $\lambda\mid \text{train}\sim \mathrm{Gamma}(a_B,b_B)$ and $y\mid\lambda\sim\mathrm{Pois}(\lambda w)$,
then the posterior predictive is negative binomial in $y$ with parameters determined by $(a_B,b_B,w)$;
for multiple points, aggregate counts and exposures and apply \eqref{eq:pp-seg}.

\paragraph{Binomial--Beta.}
If $p\mid \text{train}\sim \mathrm{Beta}(a_B,b_B)$ and $y\mid p\sim\mathrm{Binom}(m,p)$, then the
posterior predictive is Beta-Binomial:
\[
p(y\mid m,\text{train})=\binom{m}{y}\frac{B(a_B+y,b_B+m-y)}{B(a_B,b_B)}.
\]

\section{Non-conjugate GLMs: block evidence, theory, and approximations}
\label{sec:nonconj}

\paragraph{Plate model (non-conjugate single segment with PG augmentation).}
The generative graph for a non-conjugate block $(i,j]$ is identical to Figure~\ref{fig:plate-ef}
with $\theta\sim\pi(\theta)$ replaced by a non-conjugate prior (e.g., Gaussian on a log or logit
scale). For Pólya--Gamma-augmented binomial logistic blocks (Theorem~\ref{thm:pg}), introducing
the latent scale-mixture variables $\omega_t\sim\mathrm{PG}(m_t,0)$ restores conditional
conjugacy and is reflected in the plate as a second observed-indexed latent:
\begin{equation}
\theta\sim\mathcal{N}(\nu,\rho^2),\quad
\omega_t\stackrel{\mathrm{ind}}{\sim}\mathrm{PG}(m_t,0),\quad
y_t\mid\theta,\omega_t,m_t\stackrel{\mathrm{ind}}{\sim}\mathrm{Binom}(m_t,\sigma(\theta))\text{ with \eqref{eq:pg-identity}},
\label{eq:plate-nonconj}
\end{equation}
where $\sigma$ is the logistic function.
\begin{figure}[H]
\centering
\begin{tikzpicture}
  \node[obs] (y) {$y_t$};
  \node[latent, above=0.9cm of y] (om) {$\omega_t$};
  \node[latent, left=1.2cm of om] (th) {$\theta$};
  \node[const, left=0.9cm of th] (nr) {$\nu,\rho^2$};
  \node[const, right=0.9cm of y] (m) {$m_t$};
  \edge {nr} {th};
  \edge {th,om,m} {y};
  \edge {m} {om};
  \plate {p1} {(y)(om)} {$t\in(i,j]$};
\end{tikzpicture}
\caption{Plate model for a non-conjugate block with Pólya--Gamma augmentation. Conditional on $\omega_t$, the block is quadratic in $\theta$ with a Gaussian posterior (Proposition~\ref{prop:pg-mf}); the same template applies to Laplace and variational routines with the $\omega_t$ node omitted.}
\label{fig:plate-nonconj}
\end{figure}

In many applications the single-segment parameterization is not conjugate to the observation
model—for example, log-normal priors for positive means, heavy-tailed shrinkage priors, or
logistic/negative-binomial links with structural offsets. This section develops a plug-in
toolkit for computing approximate block evidences $A^{(0)}_{ij}$ and moments $A^{(r)}_{ij}$
when conjugacy is unavailable, while keeping the dynamic-programming layer of
Section~\ref{sec:dp} unchanged. When approximate block evidences are used, the DP is exact for
the approximate block-score model; additional error statements require explicit blockwise
approximation control.

We focus on one-dimensional (or low-dimensional) canonical GLMs, where for a block $(i,j]$ we
must approximate integrals of the form
\[
A^{(0)}_{ij}
 = \int \Big[\prod_{t=i+1}^j p(y_t\mid\theta)\Big] p(\theta)\,d\theta,\qquad
A^{(r)}_{ij}
 = \int \Big[\prod_{t=i+1}^j p(y_t\mid\theta)\Big] p(\theta)\,m(\theta)^r\,d\theta,
\]
with $m(\theta):=\mathbb{E}[Y\mid\theta]$. We describe four complementary approaches:
(i) first-order Laplace approximation; (ii) variational lower bounds (Jaakkola–Jordan \citep{jaakkola2000logisticvb} type
for logistic regression and mean-field for more general GLMs); (iii) expectation propagation
(EP); and (iv) P\'olya--Gamma augmentation for logistic and negative-binomial blocks. Throughout
we retain weights $w_t$ and cumulative statistics $(S_{ij},W_{ij},H_{ij})$.

\subsection{Model and notation for GLM blocks}

Consider a one-parameter canonical GLM on a block $(i,j]$:
\[
\log p(y_t\mid\theta)
 = w_t\big\{ y_t\theta - b(\theta)\big\} + \log h(y_t,w_t),
\]
where $b$ is the cumulant function and $h$ collects base measures (including any $w_t$-dependent
terms). The block log-likelihood is
\[
\ell_{ij}(\theta)
 := \sum_{t=i+1}^j w_t\{y_t\theta - b(\theta)\} + H_{ij},
\qquad
H_{ij}:=\sum_{t=i+1}^j \log h(y_t,w_t).
\]
Let $\pi(\theta)=p(\theta)$ be an arbitrary prior (not necessarily conjugate), and write
$\Psi_{ij}(\theta):=\ell_{ij}(\theta)+\log\pi(\theta)$ for the block log-posterior.

Derivatives are
\[
\ell_{ij}'(\theta) = S_{ij} - W_{ij} b'(\theta),\qquad
\ell_{ij}''(\theta) = -W_{ij} b''(\theta),
\]
and
\[
\Psi_{ij}'(\theta) = \ell_{ij}'(\theta) + (\log\pi)'(\theta),\qquad
\Psi_{ij}''(\theta) = \ell_{ij}''(\theta) + (\log\pi)''(\theta).
\]

\subsection{Existence and uniqueness of block modes}

\begin{lemma}[Strict log-concavity and uniqueness]
\label{lem:concave-unique}
Suppose $b''(\theta)>0$ for all $\theta$ (canonical GLM), the prior is log-concave
$((\log\pi)''(\theta)\le 0)$, and $W_{ij}>0$. Then $\Psi_{ij}$ is strictly concave. If, in
addition, $\Psi_{ij}$ is proper and attains its supremum on the parameter domain (for example by
coercive tails or compact-domain truncation), the maximizer $\hat\theta_{ij}$ exists and is unique.
\end{lemma}

\begin{proof}
Since $W_{ij}>0$ and $b''(\theta)>0$, we have $-W_{ij}b''(\theta)<0$. Adding
$(\log\pi)''(\theta)\le 0$ yields $\Psi_{ij}''(\theta)<0$ wherever the derivatives exist, so
$\Psi_{ij}$ is strictly concave. Strict concavity gives uniqueness of any maximizer; existence is
provided by the stated attainment assumption.
\end{proof}

\subsection{Laplace approximation for single-block evidence}

\begin{theorem}[First-order Laplace approximation]
\label{thm:laplace}
Let $a\asymp b$ mean that $a/b$ is bounded above and below by positive constants independent of
$(i,j)$. Assume:
\begin{enumerate}
\item[(A1)] $\Psi_{ij}$ is three times continuously differentiable on a neighborhood of its
unique interior maximizer $\hat\theta_{ij}$;
\item[(A2)] $H_{ij}^{\star}:=-\Psi_{ij}''(\hat\theta_{ij})>0$ and $H_{ij}^{\star}\asymp W_{ij}$ as $W_{ij}\to\infty$;
\item[(A3)] $|\Psi_{ij}^{(3)}(\theta)|\le C W_{ij}$ on a neighborhood of $\hat\theta_{ij}$ for a constant $C$
independent of $(i,j)$;
\item[(A4)] the contribution of the integral outside that neighborhood is negligible at the stated
Laplace order relative to the local quadratic contribution. This is an assumption on the block and
prior; log-concavity helps with uniqueness but does not by itself provide the required tail bound
without additional coercivity or compact-domain control.
\end{enumerate}
Then the block evidence admits the first-order Laplace expansion
\begin{equation}
A^{(0)}_{ij}
= \exp\{\Psi_{ij}(\hat\theta_{ij})\}\,\Big(\frac{2\pi}{H_{ij}^{\star}}\Big)^{1/2}
\Big[1+\mathcal{O}(W_{ij}^{-1})\Big],
\label{eq:laplace-A0-mult}
\end{equation}
and therefore
\begin{equation}
\log A^{(0)}_{ij}
= \Psi_{ij}(\hat\theta_{ij}) + \tfrac{1}{2}\log(2\pi)
  - \tfrac{1}{2}\log H_{ij}^{\star} + \mathcal{O}(W_{ij}^{-1}).
\label{eq:laplace-A0}
\end{equation}
Moreover, for any twice continuously differentiable test function $g$ with bounded derivatives near
$\hat\theta_{ij}$, writing $M^{[g]}_{ij}:=\int e^{\Psi_{ij}(\theta)}g(\theta)\,d\theta$ (an
unnormalized weighted integral),
\begin{equation}
\frac{M^{[g]}_{ij}}{A^{(0)}_{ij}}
= g(\hat\theta_{ij}) + \mathcal{O}(W_{ij}^{-1}).
\label{eq:laplace-moment}
\end{equation}
\end{theorem}

\begin{proof}
Write $u=\sqrt{H_{ij}^{\star}}(\theta-\hat\theta_{ij})$. A third-order Taylor expansion of
$\Psi_{ij}$ around $\hat\theta_{ij}$ gives
\[
\Psi_{ij}(\theta)
= \Psi_{ij}(\hat\theta_{ij}) - \tfrac12 u^2 + R_{ij}(u),
\qquad |R_{ij}(u)| \le C' |u|^3 / \sqrt{W_{ij}}
\]
on the local neighborhood, using Assumptions~(A2)--(A3). Integrating the leading quadratic term yields the
Gaussian factor $(2\pi/H_{ij}^{\star})^{1/2}$, while the remainder contributes a relative
$\mathcal{O}(W_{ij}^{-1})$ correction by the standard Laplace method; Assumption~(A4) controls the tails and
establishes \eqref{eq:laplace-A0-mult}--\eqref{eq:laplace-A0}. Applying the same expansion to the numerator
$M^{[g]}_{ij}$ and using a second-order Taylor expansion of $g$ around $\hat\theta_{ij}$ yields
\eqref{eq:laplace-moment}. We state the first-order ratio result: any explicit second-order
correction depends jointly on derivatives of $g$ and of $\Psi_{ij}$ and is not captured by a universal
$g''/(2H_{ij}^{\star})$ term.
\end{proof}

\paragraph{Efficient Newton updates using block summaries.}
For canonical GLMs,
\[
\Psi_{ij}'(\theta) = S_{ij} - W_{ij} b'(\theta) + (\log\pi)'(\theta),\quad
\Psi_{ij}''(\theta) = -W_{ij} b''(\theta) + (\log\pi)''(\theta),
\]
so Newton iteration is
\[
\theta^{(m+1)} \leftarrow \theta^{(m)}
 - \frac{S_{ij} - W_{ij} b'(\theta^{(m)}) + (\log\pi)'(\theta^{(m)})}
        {-\,W_{ij} b''(\theta^{(m)}) + (\log\pi)''(\theta^{(m)})},
\]
with cost independent of block length once $S_{ij},W_{ij}$ are available.


\paragraph{Implementation.}
Algorithm~\ref{alg:laplace} implements the Laplace block-evidence approximation using a Newton or quasi-Newton solver for the block mode and the corresponding Hessian correction.
The approximation is local but often highly accurate in moderate-to-large blocks.

\begin{algorithm}[H]
\caption{GLMBlock\_Laplace$(y,w,i,j,\texttt{b},\texttt{prior})$}
\label{alg:laplace}
\KwIn{$y_{i+1:j}$, weights $w_{i+1:j}$; cumulant $b(\theta)$; prior $\pi(\theta)$}
Compute $S\leftarrow \sum_{t=i+1}^j w_t y_t$;\quad $W\leftarrow \sum_{t=i+1}^j w_t$;\quad $H\leftarrow \sum_{t=i+1}^j \log h(y_t,w_t)$\;
Define $\Psi(\theta)=H + S\theta - W b(\theta) + \log \pi(\theta)$ and its derivatives\;
Initialize $\theta^{(0)}$ (e.g., solve $b'(\theta)=S/W$ or use prior mean)\;
\For{$m=0,1,\dots$ until convergence}{
  $g\leftarrow S - W b'(\theta^{(m)}) + (\log\pi)'(\theta^{(m)})$\;
  $h\leftarrow - W b''(\theta^{(m)}) + (\log\pi)''(\theta^{(m)})$\;
  \uIf{$|g|/|h|<\varepsilon_{\text{newton}}$}{break}
  $\theta^{(m+1)} \leftarrow \theta^{(m)} - g/h$\;
}
$\hat\theta\leftarrow \theta^{(m)}$;\quad $H^\star\leftarrow -\,\Psi''(\hat\theta)$\;
$\log \widehat{A}^{(0)}_{ij}\leftarrow \Psi(\hat\theta) + \tfrac{1}{2}\log(2\pi) - \tfrac{1}{2}\log H^\star$\;
\KwOut{$\widehat{A}^{(0)}_{ij}$ and (optional) moments via \eqref{eq:laplace-moment}}
\end{algorithm}


\subsection{Variational lower bounds for GLM blocks}

When a deterministic lower bound is desirable, introduce a variational distribution $q(\theta)$
and apply Jensen:
\[
\log A^{(0)}_{ij}
=\log \int q(\theta)\,\frac{e^{\Psi_{ij}(\theta)}}{q(\theta)}\,d\theta
\ \ge\
\int q(\theta)\,\Psi_{ij}(\theta)\,d\theta + \mathcal{H}[q]
 \;:=\; \mathcal{L}_{ij}(q).
\]

\emph{Jaakkola--Jordan bound for logistic blocks.}
For Bernoulli/logistic likelihoods, a standard quadratic bound yields a tractable Gaussian
variational posterior. One convenient form (for sigmoid $\sigma$) is:
\[
\log\sigma(z)\ \ge\ \log\sigma(\xi)+\frac{z-\xi}{2}-\lambda(\xi)(z^2-\xi^2),
\qquad
\lambda(\xi)=\frac{\tanh(\xi/2)}{4\xi},
\]
with variational parameter $\xi>0$. Applying this bound to the appropriate signed linear
predictors yields a quadratic surrogate in $\theta$, so the optimal Gaussian $q(\theta)$ and
$\xi$ admit closed-form coordinate updates.

\begin{proposition}[Monotone improvement of variational block bounds]
\label{prop:var-monotone}
Coordinate-ascent updates on variational parameters (e.g.\ $\xi$ for logistic blocks) and on
a Gaussian $q(\theta)$ increase $\mathcal{L}_{ij}(q)$ until a stationary point; the resulting
$\underline{A}^{(0)}_{ij}:=\exp\{\mathcal{L}_{ij}(q)\}$ is a certified lower bound on
$A^{(0)}_{ij}$.
\end{proposition}

\begin{proof}
Each coordinate update is an exact maximization of $\mathcal{L}_{ij}(q)$ with respect to a
subset of variables holding others fixed. Therefore $\mathcal{L}_{ij}(q)$ is non-decreasing
under the updates. Because $\mathcal{L}_{ij}$ is a Jensen lower bound on $\log A^{(0)}_{ij}$,
exponentiating yields a valid evidence lower bound.
\end{proof}


\paragraph{Implementation.}
Algorithm~\ref{alg:jj} provides pseudocode for computing the Jaakkola--Jordan (JJ) variational lower bound and its associated fixed-point updates at the block level.

\begin{algorithm}[H]
\caption{LogisticBlock\_JJ$(y,m,w,i,j,\texttt{prior})$ \quad (variational lower bound)}
\label{alg:jj}
\KwIn{Binomial counts $y_t\sim\mathrm{Binom}(m_t,p)$, logit link; prior $\pi(\theta)$}
Initialize $\xi$ and a Gaussian $q(\theta)=\mathcal{N}(\mu_q,\sigma_q^2)$\;
\Repeat{converged}{
  Update the quadratic bound parameter $\xi\leftarrow \sqrt{\mathbb{E}_q[\theta^2]}$\;
  Update Gaussian $q(\theta)$ by maximizing the block ELBO $\mathcal{L}_{ij}(q)$ (closed-form for quadratic surrogate)\;
}
Return $\underline{A}^{(0)}_{ij}=\exp\{\mathcal{L}_{ij}(q)\}$ and moments from $q$.
\end{algorithm}


\subsection{Expectation propagation (EP) for GLM blocks}
Expectation propagation (EP) \citep{minka2001ep} provides a flexible deterministic approximation to the block posterior and block evidence by iteratively refining local site approximations and matching moments.
EP approximates each likelihood term $p(y_t\mid\theta)$ by a Gaussian site
$\tilde t_t(\theta)\propto\exp\{-\tfrac{1}{2}a_t\theta^2+b_t\theta\}$ so that moments match those
of the tilted distribution. After convergence, the approximate posterior is Gaussian and yields
an approximate evidence $\tilde A^{(0)}_{ij}=\int \pi(\theta)\prod \tilde t_t(\theta)\,d\theta$.
EP is neither a bound nor guaranteed to converge, but is often accurate in canonical GLMs.


\paragraph{Implementation.}
Algorithm~\ref{alg:ep} summarizes an expectation-propagation (EP) routine for approximating the GLM block evidence, highlighting the cavity/site updates and moment matching.

\begin{algorithm}[H]
\caption{GLMBlock\_EP$(y,w,i,j,\texttt{prior})$}
\label{alg:ep}
\KwIn{Canonical GLM; prior $\pi(\theta)$}
Initialize site parameters $\{a_t,b_t\}_{t=i+1}^j$; approximate posterior $\tilde p(\theta)\propto\pi(\theta)\prod \tilde t_t(\theta)$\;
\Repeat{moments stable}{
  \For{$t=i+1$ \KwTo $j$}{
    Form cavity $\tilde p^{\setminus t}(\theta)\propto \tilde p(\theta)/\tilde t_t(\theta)$\;
    Form tilted $p^{\text{tilt}}(\theta)\propto \tilde p^{\setminus t}(\theta)\,p(y_t\mid\theta)$\;
    Moment-match $p^{\text{tilt}}$ and update site $\tilde t_t(\theta)$\;
  }
}
Return approximate evidence $\tilde A^{(0)}_{ij}=\int \tilde p(\theta)\,d\theta$ and moments from $\tilde p$.
\end{algorithm}


\subsection{Pólya–Gamma augmentation (logistic and negative-binomial)}
\begin{theorem}[P\'olya--Gamma identity and conditional conjugacy]
\label{thm:pg}
For binomial logistic blocks with counts $m_t$ and successes $y_t$, let $w_t\ge 0$ denote either
one or a declared likelihood-power/replication weight, define
$b_t=w_t m_t$ and $\kappa_t=w_t(y_t-\tfrac{m_t}{2})$, and introduce latent
$\omega_t\sim\mathrm{PG}(b_t,0)$. Then the weighted logistic likelihood term admits the Gaussian
scale-mixture representation
\begin{equation}
\left\{\frac{(e^\theta)^{y_t}}{(1+e^\theta)^{m_t}}\right\}^{w_t}
 = 2^{-b_t} e^{\kappa_t\theta}
   \int_0^\infty e^{-\omega_t\theta^2/2}\,p(\omega_t)\,d\omega_t.
\label{eq:pg-identity}
\end{equation}
Conditionally on $\omega_{i+1:j}$, the block likelihood is quadratic in $\theta$:
\begin{equation}
\log p(y_{i+1:j}\mid \theta,\omega)
 =
 -\tfrac{1}{2}\theta^2 \sum_{t=i+1}^j \omega_t
 + \theta \sum_{t=i+1}^j \kappa_t
 + \mathrm{const}.
\label{eq:pg-cond-quad}
\end{equation}
With a Gaussian prior $\pi(\theta)=\mathcal{N}(\nu,\rho^2)$, the conditional posterior
$p(\theta\mid y,\omega)$ is Gaussian.
\end{theorem}

\begin{proof}
\ProofStep{Step 1: the identity.}
Equation \eqref{eq:pg-identity} is the defining P\'olya--Gamma augmentation of
\citet{polson2013pg}: it expresses the binomial-logistic term as a Gaussian scale mixture in
$\theta$.

\ProofStep{Step 2: condition on $\omega$ and collect quadratic terms.}
Given $\omega_t$, each weighted term contributes $\kappa_t\theta-\tfrac12\omega_t\theta^2$ up to
constants, where the weight has already been absorbed into $b_t$ and $\kappa_t$. Summing over $t$
yields \eqref{eq:pg-cond-quad}.

\ProofStep{Step 3: combine with a Gaussian prior.}
Multiplying the quadratic-in-$\theta$ conditional likelihood by a Gaussian prior produces a
Gaussian conditional posterior, because Gaussian priors are conjugate to Gaussian likelihoods.
\end{proof}

\begin{proposition}[Mean-field variational approximation under P\'olya--Gamma]
\label{prop:pg-mf}
Under the same setup as Theorem~\ref{thm:pg} and a Gaussian prior $\pi(\theta)=\mathcal{N}(\nu,\rho^2)$,
the mean-field factorization $q(\theta,\omega)=q(\theta)\prod_{t=i+1}^j q(\omega_t)$ admits the
following closed-form coordinate-ascent updates:
(i) $q(\theta)=\mathcal{N}(\mu_q,\sigma_q^2)$ with
$\sigma_q^{-2}=\rho^{-2}+\sum_t \mathbb{E}_q[\omega_t]$ and
$\mu_q=\sigma_q^2(\rho^{-2}\nu+\sum_t \kappa_t)$;
(ii) $q(\omega_t)=\mathrm{PG}(b_t, c_t)$ with $c_t^2=\mathbb{E}_q[\theta^2]=\mu_q^2+\sigma_q^2$,
giving $\mathbb{E}_q[\omega_t]=\tfrac{b_t}{2 c_t}\tanh(c_t/2)$. These updates are monotone in the
block ELBO.
\end{proposition}

\begin{proof}
Standard mean-field variational calculus applied to \eqref{eq:pg-cond-quad} combined with the prior.
The explicit update for $\mathbb{E}_q[\omega_t]$ uses the mean of a $\mathrm{PG}(m,c)$ distribution;
see Appendix~\ref{app:pg-mf-derivation} for the full derivation and ELBO monotonicity argument.
\end{proof}


\paragraph{Implementation.}
Algorithm~\ref{alg:pg-vb} gives a Pólya--Gamma variational Bayes (PG-VB) routine for logistic (and related) blocks, which yields a tractable surrogate block evidence by alternating between Gaussian updates and latent-variable expectations.

\begin{algorithm}[H]
\caption{GLMBlock\_PG\_VB$(y,m,i,j,\nu,\rho^2)$ \quad (Gaussian prior, PG mean-field)}
\label{alg:pg-vb}
\KwIn{Binomial/logistic block with trial counts $m_t$; Gaussian prior $\mathcal{N}(\nu,\rho^2)$}
Initialize $q(\omega_t)$ and $q(\theta)=\mathcal{N}(\mu,\sigma^2)$\;
\Repeat{ELBO converged}{
  Set $b_t\leftarrow w_t m_t$ and $\kappa_t\leftarrow w_t(y_t-m_t/2)$ if likelihood-power weights are used; otherwise $w_t=1$\;
  $\sigma^{-2}\leftarrow \rho^{-2} + \sum_{t=i+1}^j \mathbb{E}_q[\omega_t]$\;
  $\mu \leftarrow \sigma^2 \left(\rho^{-2}\nu + \sum_{t=i+1}^j \kappa_t\right)$\;
  \ForEach{$t$}{
    Update $q(\omega_t)=\mathrm{PG}(w_t m_t,c_t)$ using current $\mathbb{E}_q[\theta^2]$\;
  }
}
Return a lower bound $\underline{A}^{(0)}_{ij}$ and approximate moments from $q(\theta)$.
\end{algorithm}


\paragraph{Using approximate block evidences in BayesBreak.}
Once any block-evidence approximation is available (Laplace, JJ, EP, PG-VB, or another surrogate), BayesBreak uses it exactly as an input score array: build a triangular array of approximate block log-evidences and run the same DP recursions. The resulting posterior summaries are exact for that approximate score array, not automatically exact for the original non-conjugate model.
Algorithm~\ref{alg:precompute-nonconj} provides a generic recipe for constructing this array in a model-agnostic way, emphasizing reuse of block-local optimizers and cached sufficient-statistic summaries.

\begin{algorithm}[H]
\caption{Precompute block evidences for non-conjugate GLMs}
\label{alg:precompute-nonconj}
\KwIn{$\{(x_t,y_t,w_t)\}_{t=1}^n$, method $\in\{$Laplace, JJ, PG--VB, EP$\}$}
\For{$i=0$ \KwTo $n-1$}{
  Initialize running summaries (as needed)\;
  \For{$j=i+1$ \KwTo $n$}{
    Update block summaries for $(i,j]$\;
    Call chosen block routine to compute $\widehat{A}^{(0)}_{ij}$ (and optional moments)\;
  }
}
Run DP (Algorithm~\ref{alg:dp-core}) with resulting block evidences.
\end{algorithm}


\subsection{Propagation of block approximation error through the DP}

\begin{proposition}[Log-evidence and posterior-odds stability]
\label{prop:stability}
Let $\widehat{A}^{(0)}_{ij}$ be approximate block evidences with log-errors
$\delta_{ij}:=\log \widehat{A}^{(0)}_{ij}-\log A^{(0)}_{ij}$. Assume the \emph{uniform blockwise
bound}
\begin{equation}
\begin{aligned}
|\delta_{ij}|\le \varepsilon \ &\text{uniformly over all candidate blocks $(i,j]$, $0\le i<j\le n$,}\\
&\text{reachable by some $k$-segmentation with $k\le k_{\max}$.}
\end{aligned}
\label{eq:stability-uniform}
\end{equation}
Then for every DP state $(k,j)$ with $k\le k_{\max}$,
\[
e^{-k\varepsilon}L_{k j}\le \widehat{L}_{k j}\le e^{k\varepsilon}L_{k j},
\qquad
e^{-k\varepsilon}R_{k j}\le \widehat{R}_{k j}\le e^{k\varepsilon}R_{k j}.
\]
Consequently:
\begin{enumerate}
\item For any $k,k'\le k_{\max}$, the posterior odds for the segment count satisfy
\begin{equation}
\left|
\log\frac{\widehat{P}(k\mid y)}{\widehat{P}(k'\mid y)}
-\log\frac{P(k\mid y)}{P(k'\mid y)}
\right|
\le (k+k')\varepsilon.
\label{eq:stability-k-odds}
\end{equation}
\item For any fixed $k\le k_{\max}$ and any two candidate locations $h,h'$ of boundary $t_p$,
\begin{equation}
\left|
\log\frac{\widehat{P}(t_p=h\mid y,k)}{\widehat{P}(t_p=h'\mid y,k)}
-\log\frac{P(t_p=h\mid y,k)}{P(t_p=h'\mid y,k)}
\right|
\le 2k\varepsilon.
\label{eq:stability-b-odds}
\end{equation}
\end{enumerate}
Thus uniform blockwise log-evidence control yields uniform control of global posterior \emph{odds}. Turning
these odds bounds into absolute probability error bounds requires an additional margin assumption and is
therefore left implicit.
\end{proposition}

\begin{proof}
Any term contributing to $L_{k j}$ is a product of exactly $k$ block evidences. Hence the corresponding
approximate term differs by a multiplicative factor in $[e^{-k\varepsilon},e^{k\varepsilon}]$. Summing over all
paths yields the same sandwich bound for $\widehat{L}_{k j}$, and the suffix bound is identical. Equation
\eqref{eq:stability-k-odds} follows immediately because $P(k\mid y)$ is proportional to $p(k)L_{k n}/C_k$ and the
prior and normalization constants are unchanged. For boundary odds, note that
$P(t_p=h\mid y,k)\propto L_{p,h}R_{k-p,h}$. The numerator for each candidate location involves $k$ block factors in
total, so each candidate-specific weight is perturbed by at most $e^{\pm k\varepsilon}$; comparing two candidates
therefore yields the factor $e^{\pm 2k\varepsilon}$ and hence \eqref{eq:stability-b-odds}.
\end{proof}

\begin{corollary}[Preservation of boundary rankings under a margin condition]
\label{cor:ranking}
Under the hypotheses of Proposition~\ref{prop:stability}, if there is a \emph{log-odds margin}
$\Delta>0$ separating the best candidate $h^\star$ from every alternative,
\[
\log\frac{P(t_p=h^\star\mid y,k)}{P(t_p=h\mid y,k)}\ge \Delta
\quad\text{for all }h\ne h^\star,
\]
then for every $\varepsilon<\Delta/(2k)$ the approximate posterior preserves the exact modal
boundary location: $\arg\max_h \widehat P(t_p=h\mid y,k)=h^\star$.
\end{corollary}

\begin{proof}
By \eqref{eq:stability-b-odds}, the log-odds between $h^\star$ and any alternative $h$ is
perturbed by at most $2k\varepsilon<\Delta$, so the exact log-odds positivity is preserved.
\end{proof}

\begin{remark}[A small numerical illustration of the stability bound]
\label{rem:stability-numerical}
Consider $n=500$, $k_{\max}=10$, and a non-conjugate block routine whose verified uniform
per-block log-evidence error over all reachable blocks is $\varepsilon=0.02$. Proposition~\ref{prop:stability}
then guarantees segment-count log-odds error at most $(5+10)\times0.02=0.30$ when comparing
$k=5$ to $k=10$, and boundary log-odds error at most $2\cdot10\cdot0.02=0.40$ under any fixed
$k\le10$. In terms of probability ratios, these correspond to factors of about $e^{0.30}$ and
$e^{0.40}$. The numbers in this illustration are conditional on having already verified the
uniform block-error bound; they are not a generic guarantee for any chosen approximation routine.
\end{remark}

\begin{remark}[When each non-conjugate routine struggles]
\label{rem:failure-modes}
The stability bound in Proposition~\ref{prop:stability} is tight only when the uniform blockwise
error $\varepsilon$ is actually small. Each non-conjugate routine has its own failure mode:
Laplace (Theorem~\ref{thm:laplace}) requires posterior concentration near $\hat\theta_{ij}$, so
it can be inaccurate on very short segments ($W_{ij}$ small) or for heavy-tailed priors;
mean-field variational Bayes (Proposition~\ref{prop:pg-mf}) can be overconfident and
underestimate posterior variance, biasing moment numerators; EP can fail to converge for
non-log-concave likelihoods; 1D quadrature is accurate for smooth 1D integrands but expensive
in higher dimensions. For these reasons, monitor $\varepsilon$ empirically per
dataset rather than assuming it. A practical validation checklist is to compare a subset of blocks against
quadrature or another trusted reference, report both maximum and high-quantile block errors, inspect
errors along the exported MAP path, and assess sensitivity of $P(k\mid y)$ and boundary marginals.
\end{remark}

\paragraph{Approximation-validation checklist.}
Before interpreting non-conjugate posterior summaries, compare a subset of block evidences against
higher-accuracy quadrature or another trusted reference; report maximum and high-quantile block
errors over reachable blocks; inspect errors on the MAP path; and assess sensitivity of $p(k\mid y)$,
boundary marginals, and exported segmentations to the approximation method. The ranking corollary
is meaningful only when the exact posterior margin is larger than the propagated approximation error.


\section{Prediction: group likelihoods and MAP/Bayes signal evaluation}
\label{sec:prediction}

Having developed training-time posterior inference for partitions and segment parameters (Sections~\ref{sec:dp}--\ref{sec:nonconj}), we now turn to prediction for new $(X,y)$ data.
The prediction layer consumes posterior summaries (e.g., boundary posteriors, segment-parameter posteriors, and group evidences) and is therefore agnostic to whether block evidences were computed in closed form (Section~\ref{sec:families}) or via approximation (Section~\ref{sec:nonconj}). We assume we have already trained (or otherwise obtained)
a collection of $G$ group-specific BayesBreak models
$\{\mathcal{M}_g\}_{g=1}^G$, where each $\mathcal{M}_g$ comprises:
(i) a fitted segmentation posterior (e.g.\ $P_g(k\mid y)$ and boundary posteriors),
(ii) an exported MAP segmentation $\widehat{t}^{(g)}$ with segment-level posterior
hyperparameters, and (iii) optionally a Bayesian curve representation
(e.g.\ posterior mean/variance at design points).

Typical outputs are group log scores, posterior group probabilities, optional unit-level
responsibilities, exported MAP signal values, Bayes-curve values, and optional pointwise
uncertainty summaries. These prediction interfaces define the outputs populated by finalized diagnostics for the placeholder real-data examples.

Prediction supports two primary tasks:

\paragraph{(P1) Group membership likelihoods/posteriors.}
Given new inputs $(X,y)$, compute for each group $g$ an out-of-sample
posterior-predictive likelihood
$p(\text{new}\mid \mathcal{M}_g)$ (or its log), and convert these into posterior group
membership probabilities $P(g\mid \text{new})$ under a prior $\pi_g$.

\paragraph{(P2) Signal prediction given group.}
Given query design points $X$ and a specified (or inferred) group label $g$, return:
(i) MAP piecewise-constant signal values $\widehat{f}^{\text{MAP}}_g(X)$ including the
corresponding breakpoints and segment levels, and/or (ii) Bayesian curve values
$\widehat{f}^{\text{Bayes}}_g(X)$ with pointwise uncertainty summaries when available.

\begin{table}[H]
\centering
\small
\begin{tabular}{@{}ll@{}}
\toprule
Prediction output & Interpretation \\
\midrule
$\ell_g$ & Group-specific log posterior-predictive score. \\
$P(g\mid \mathrm{new})$ & Group posterior after combining $\ell_g$ with prior weight $\pi_g$. \\
$\widehat f^{\mathrm{MAP}}_g(X)$ & Piecewise-constant signal under the exported MAP segmentation. \\
$\widehat f^{\mathrm{Bayes}}_g(X)$ & Bayes-curve posterior mean, optionally with pointwise uncertainty. \\
Unit responsibilities & Optional probabilities for set-valued units or grouped observations. \\
\bottomrule
\end{tabular}
\caption[Prediction-layer outputs]{Prediction-layer outputs exposed by fitted BayesBreak models. These interfaces support downstream scoring and calibration separately from the empirical result tables.}
\label{tab:prediction-outputs}
\end{table}

Crucially, we distinguish \emph{single point observations} $(x_i,y_i)$ from
\emph{multivariate / set-valued observations} where each unit contains multiple internal
pairs $(x_{u r},y_{u r})$, i.e.\ $X_u$ and $y_u$ are vectors. We also separately address
\emph{vector-valued responses} $y_i\in\mathbb{R}^d$ observed at each design point.

\subsection{Prediction inputs: pointwise vs multivariate/set-valued observations}
\label{sec:prediction-inputs}

\paragraph{Case A (pointwise observations).}
We observe a single ordered sequence
\[
\mathcal{D}^{\text{new}}_{\text{pt}}
:= \{(x_i^{\text{new}},y_i^{\text{new}},w_i^{\text{new}})\}_{i=1}^{m},
\qquad x_1^{\text{new}}<\cdots<x_m^{\text{new}},
\]
with optional exposure or precision weights $w_i^{\text{new}}$. The default is
$w_i^{\text{new}}=1$; use inter-point gaps only when they are genuine exposures or
known-precision weights for the observation model, not merely because the design is irregular.

\paragraph{Case B (multivariate/set-valued observations).}
We observe $U$ units (e.g.\ ``fragments'', ``windows'', or generic composite observations),
each with an interval support and multiple internal pairs:
\[
\mathcal{D}^{\text{new}}_{\text{mv}}
:= \Big\{\big([a_u,b_u],\{(x_{u r},y_{u r},w_{u r})\}_{r=1}^{R_u}\big)\Big\}_{u=1}^{U},
\qquad a_u < x_{u r}\le b_u.
\]
Here $X_u=(x_{u1},\dots,x_{uR_u})$ and $y_u=(y_{u1},\dots,y_{uR_u})$ are \emph{multivariate}
(set-valued) per-unit inputs. Prediction must score each unit by integrating or aggregating
over its internal observations, and optionally return unit-level group responsibilities.

\paragraph{Case C (vector-valued response at each design point).}
At each design point we observe a vector response $\mathbf{y}_i\in\mathbb{R}^d$ (or mixed
discrete/continuous components). We consider two standard constructions:
(i) \emph{factorized EF} across dimensions, or (ii) a \emph{multivariate EF} with a joint
log-partition $A(\theta)$ and conjugate prior (when available). The factorized case is
implementation-friendly and preserves the DP structure with localized changes, while the multivariate
case requires a genuine joint likelihood and conjugate or approximate block routine.

\subsection{Posterior-predictive scoring for a fixed segment under EF conjugacy}
\label{sec:predictive-scoring}

For group membership inference, we require out-of-sample likelihoods. The most
computationally efficient and interpretable approach is to score new data against the
\emph{exported} group curve (MAP segmentation or Bayesian curve). We first derive
segment-level posterior-predictive likelihoods under EF conjugacy; these become building
blocks for both Case A and Case B.

Fix a group $g$ and consider a segment (or block) $B$ of its exported MAP segmentation.
Let the posterior for the segment parameter $\theta$ after training be the conjugate posterior
\[
p(\theta\mid \alpha^{(g)}_B,\beta^{(g)}_B)
\ \propto\
\exp\{\eta(\theta)^\top\alpha^{(g)}_B - \beta^{(g)}_B A(\theta)\},
\]
where $(\alpha^{(g)}_B,\beta^{(g)}_B)$ are the training-updated hyperparameters for that segment.
Now let $\mathcal{Y}_B^{\text{new}}$ denote the collection of new observations whose design
points fall into this segment.

\begin{proposition}[EF posterior-predictive likelihood within a fixed segment]
\label{prop:ef-predictive}
Assume the weighted EF likelihood and conjugate prior conventions of \S\ref{sec:ef} and
\S\ref{sec:notation}, and let $(S_B^{\text{new}},W_B^{\text{new}},H_B^{\text{new}})$ denote the
EF block summaries of the new data restricted to segment $B$:
\[
S_B^{\text{new}}:=\sum_{i\in B} w_i^{\text{new}} T(y_i^{\text{new}}),\qquad
W_B^{\text{new}}:=\sum_{i\in B} w_i^{\text{new}},\qquad
H_B^{\text{new}}:=\sum_{i\in B} \log h(y_i^{\text{new}},w_i^{\text{new}}).
\]
Then the posterior-predictive likelihood of the new data in segment $B$ under group $g$ is
\begin{equation}
p(\mathcal{Y}_B^{\text{new}} \mid \mathcal{M}_g)
=
\exp\{H_B^{\text{new}}\}\,
\frac{Z(\alpha_B^{(g)}+S_B^{\text{new}},\beta_B^{(g)}+W_B^{\text{new}})}
     {Z(\alpha_B^{(g)},\beta_B^{(g)})}.
\label{eq:pp-seg}
\end{equation}
\end{proposition}

\begin{proof}
The proof is identical to the single-block evidence derivation (Theorem~\ref{thm:ef-integral})
after replacing the prior hyperparameters $(\alpha_0,\beta_0)$ by the segment posterior
hyperparameters $(\alpha_B^{(g)},\beta_B^{(g)})$ learned from training data.
\end{proof}

\paragraph{Interpretation.}
Equation~\eqref{eq:pp-seg} integrates out the segment parameter $\theta$ under the
\emph{posterior} induced by training, yielding a strictly Bayesian out-of-sample likelihood
for new observations assigned to that segment. It is therefore suitable for group scoring
and group posterior computation.

\subsection{Group likelihoods for pointwise observations (Case A)}
\label{sec:group-lik-point}

\paragraph{MAP-segmentation-based scoring.}
Let $\widehat{t}^{(g)}$ denote the exported MAP boundaries for group $g$ on the relevant domain.
Assign each new design point $x_i^{\text{new}}$ to its containing MAP segment $B=B(i;g)$ and
compute segment-wise block summaries $(S_B^{\text{new}},W_B^{\text{new}},H_B^{\text{new}})$.
The group log-likelihood is then:
\begin{equation}
\ell_g^{\text{MAP}}(\mathcal{D}^{\text{new}}_{\text{pt}})
:= \log p(\mathcal{D}^{\text{new}}_{\text{pt}}\mid \mathcal{M}_g)
= \sum_{B\in \widehat{t}^{(g)}} \log p(\mathcal{Y}_B^{\text{new}}\mid \mathcal{M}_g),
\label{eq:group-lik-map}
\end{equation}
where each segment term uses \eqref{eq:pp-seg}. This yields $\mathcal{O}(m)$ evaluation time
once segment membership is determined.

\paragraph{Bayesian-curve-based scoring (optional).}
If $\mathcal{M}_g$ exports a Bayesian curve representation that provides a local posterior over
the mean parameter at each design point (or each small bin), then one can score each point
under its local posterior predictive and sum over points. This can improve robustness near
uncertain boundaries at additional computational cost.

\paragraph{Resegmentation scoring (optional, flexible).}
An alternative definition of group likelihood is to rerun the DP on the new data under
group-specific hyperparameters (treating the new dataset as another draw from group $g$’s
prior). This yields an evidence $P(\mathcal{D}^{\text{new}}_{\text{pt}}\mid g)$ that does not
condition on the trained MAP boundaries. This option is useful when groups differ primarily in
hyperparameters rather than in the learned curve geometry, but is computationally
$\mathcal{O}(k_{\max}m^2)$ per group.

\subsection{Group likelihoods for multivariate/set-valued observations (Case B)}
\label{sec:group-lik-mv}

In Case B, each unit $u$ contains multiple internal observations
$\{(x_{u r},y_{u r})\}_{r=1}^{R_u}$. Two output granularities are typically desired:
(i) \emph{unit-level} group likelihoods and group responsibilities, and (ii) a \emph{global}
group likelihood for the entire new dataset (product over units).

\paragraph{Unit-level predictive likelihood under a fixed group.}
Fix a group $g$. Partition the unit’s internal design points by the MAP segments of group $g$.
For each segment $B$, compute the block summaries $(S_{u,B}^{\text{new}},W_{u,B}^{\text{new}},H_{u,B}^{\text{new}})$
restricted to those internal points. Then define the unit likelihood as:
\begin{equation}
\ell_{u g}
:= \log p(\text{unit }u \mid \mathcal{M}_g)
= \sum_{B \in \widehat{t}^{(g)}} \log p(\mathcal{Y}^{\text{new}}_{u,B}\mid \mathcal{M}_g),
\label{eq:unit-lik}
\end{equation}
where $\mathcal{Y}^{\text{new}}_{u,B}$ is the subset of unit $u$’s internal observations that
fall into segment $B$ and each segment term uses \eqref{eq:pp-seg} with summaries for that subset.

\paragraph{Dataset likelihood and independence assumption.}
Assuming conditional independence of units given group membership, the dataset likelihood under
group $g$ is:
\[
\ell_g^{\text{MV}}(\mathcal{D}^{\text{new}}_{\text{mv}})=\sum_{u=1}^{U}\ell_{u g}.
\]
This is the direct multivariate generalization of \eqref{eq:group-lik-map} and corresponds to
the earlier ``fragment-level'' intuition, but is fully application-agnostic.

\subsection{Vector-valued responses (Case C): factorized and multivariate EF scoring}
\label{sec:prediction-multivariate-response}

Suppose each observation is $\mathbf{y}_i=(y_{i1},\dots,y_{id})$.

\paragraph{Factorized EF across dimensions (recommended default).}
Assume conditional independence across dimensions given segment parameters:
\[
p(\mathbf{y}_i\mid \boldsymbol{\theta})
=\prod_{\ell=1}^d p_\ell(y_{i\ell}\mid \theta_\ell),
\qquad
p(\boldsymbol{\theta})=\prod_{\ell=1}^d p_\ell(\theta_\ell).
\]
Then all block evidences and posterior-predictive terms factorize over $\ell$, so
log-likelihoods are sums over dimensions:
\[
\ell_g(\mathbf{y})=\sum_{\ell=1}^d \ell_{g,\ell}(y_{\cdot\ell}),
\]
where each $\ell_{g,\ell}$ is computed by \eqref{eq:group-lik-map} or \eqref{eq:unit-lik} for
dimension $\ell$. Computationally, this costs $\mathcal{O}(d)$ per observation (or per unit).

\paragraph{Multivariate EF (when available).}
If $(\mathbf{y}_i)$ admit a multivariate EF with conjugate prior on a joint parameter
$\theta\in\mathbb{R}^p$, then the same ratio-of-normalizers structure applies with vector-valued
sufficient statistics. BayesBreak’s DP layer remains unchanged; only the definition of $S_{ij}$
and $Z(\alpha,\beta)$ changes.

\subsection{From group likelihoods to posterior group membership probabilities}
\label{sec:group-posterior}

\paragraph{Single-dataset membership (one group per dataset).}
If the new dataset is assumed to come from a single (unknown) group $g$, define:
\[
\ell_g := \log p(\text{new}\mid \mathcal{M}_g),
\qquad
P(g\mid \text{new})
= \frac{\pi_g\exp(\ell_g)}{\sum_{g'=1}^G \pi_{g'}\exp(\ell_{g'})}.
\]
This yields a \emph{group membership likelihood vector} (the $\ell_g$) and a
\emph{group membership posterior vector} (the $P(g\mid\text{new})$), matching the requested
prediction behavior in an application-agnostic way.

\paragraph{Unit-wise membership (one group per unit).}
If each multivariate unit $u$ is assumed to originate from one group, then define unit-wise
responsibilities:
\begin{equation}
r_{u g}
:= P(g\mid \text{unit }u)
= \frac{\pi_g\exp(\ell_{u g})}{\sum_{g'=1}^G \pi_{g'}\exp(\ell_{u g'})},
\label{eq:unit-resp}
\end{equation}
where $\ell_{u g}$ is given by \eqref{eq:unit-lik}. This is precisely the out-of-sample analog
of the E-step responsibilities in the latent-group model (Section~\ref{sec:latent-em}).

\subsection{MAP signal evaluation given group membership}
\label{sec:map-eval}

Given a group label $g$ and query design points $X^\star=\{x^\star_1,\dots,x^\star_M\}$,
BayesBreak returns either:
(i) \emph{MAP piecewise-constant predictions} or (ii) \emph{Bayesian curve predictions}.

\paragraph{MAP piecewise-constant predictions.}
Let $\widehat{t}^{(g)}$ be the exported MAP boundaries and let each MAP segment $B$ provide
the posterior mean of the mean parameter:
\[
\widehat{\mu}^{(g)}_B := \mathbb{E}[m(\theta)\mid \alpha^{(g)}_B,\beta^{(g)}_B].
\]
Then the MAP signal at a query point $x^\star$ is:
\[
\widehat{f}^{\text{MAP}}_g(x^\star) = \widehat{\mu}^{(g)}_{B(x^\star;g)},
\]
where $B(x^\star;g)$ denotes the MAP segment containing $x^\star$.

\paragraph{Bayesian curve predictions.}
If the DP-derived Bayes regression curve is exported on a grid, return
$\widehat{f}^{\text{Bayes}}_g(x^\star)$ by evaluating on the nearest grid point (or by
piecewise-constant interpolation within each grid interval). When pointwise posterior
variances are available, also return pointwise credible intervals; simultaneous bands require an
additional construction and are not implied by the pointwise variances alone.

\subsection{Algorithms for prediction}
\label{sec:prediction-algorithms}

\begin{algorithm}[H]
\caption{PredictGroup$(\{\mathcal{M}_g\}_{g=1}^G,\ \mathcal{D}^{\text{new}},\ \pi)$}
\label{alg:predict-group}
\KwIn{
Trained group models $\{\mathcal{M}_g\}$; new data $\mathcal{D}^{\text{new}}$ (Case A or Case B);
group prior weights $\pi$ (default uniform)
}
\KwOut{
Log-likelihoods $\{\ell_g\}$ and posterior membership probabilities $\{P(g\mid \text{new})\}$;
if Case B, also unit responsibilities $\{r_{u g}\}$
}
\BlankLine
\For{$g=1$ \KwTo $G$}{
  \eIf{\textbf{Case A (pointwise)}}
  {
    Assign each $(x_i^{\text{new}},y_i^{\text{new}})$ to a MAP segment $B(i;g)$ under $\widehat{t}^{(g)}$\;
    Accumulate segment summaries $(S_B^{\text{new}},W_B^{\text{new}},H_B^{\text{new}})$\;
    Compute $\ell_g \leftarrow \sum_B \Big[ H_B^{\text{new}} + \log Z(\alpha_B^{(g)}+S_B^{\text{new}},\beta_B^{(g)}+W_B^{\text{new}})
      - \log Z(\alpha_B^{(g)},\beta_B^{(g)})\Big]$\;
  }
  {
    \tcp{Case B (multivariate/set-valued units)}
    \For{$u=1$ \KwTo $U$}{
      Partition unit $u$ internal pairs by MAP segments of group $g$; compute per-segment summaries\;
      Compute $\ell_{u g}$ via \eqref{eq:unit-lik} using \eqref{eq:pp-seg}\;
    }
    $\ell_g \leftarrow \sum_{u=1}^U \ell_{u g}$\;
  }
}
Compute $P(g\mid \text{new}) \propto \pi_g\exp(\ell_g)$ and normalize over $g$\;
\If{\textbf{Case B}}{
  Compute unit responsibilities $r_{u g}\propto \pi_g\exp(\ell_{u g})$ and normalize over $g$ for each $u$ (cf.\ \eqref{eq:unit-resp})\;
}
\end{algorithm}

\begin{algorithm}[H]
\caption{PredictMAPSignal$(\mathcal{M}_g,\ X^\star,\ \texttt{mode})$}
\label{alg:predict-map}
\KwIn{
A group model $\mathcal{M}_g$ with exported MAP segmentation $\widehat{t}^{(g)}$ and/or Bayes curve;
query points $X^\star=\{x^\star_1,\dots,x^\star_M\}$; mode $\in\{\texttt{MAP},\texttt{Bayes}\}$
}
\KwOut{
Predicted signal values at $X^\star$ (and optionally uncertainty)
}
\BlankLine
\eIf{\texttt{mode}=\texttt{MAP}}{
  \ForEach{$x^\star$ in $X^\star$}{
    Find containing MAP segment $B(x^\star;g)$ (binary search on breakpoint coordinates)\;
    Return $\widehat{f}^{\text{MAP}}_g(x^\star)=\widehat{\mu}^{(g)}_{B(x^\star;g)}$\;
  }
}{
  Evaluate/interpolate the exported Bayes regression curve at $X^\star$; optionally return posterior variances\;
}
\end{algorithm}

\subsection{Complexity}
For Case A, prediction under a fixed group is $\mathcal{O}(m)$ after segment assignment
($\mathcal{O}(m\log k)$ worst case via binary search; often $\mathcal{O}(m)$ with a streaming pointer).
For Case B, it is $\mathcal{O}(\sum_u R_u)$ per group, plus segment assignment for internal points.
Vector-valued responses add a factor of $d$ in the factorized EF construction. Importantly,
prediction does \emph{not} require rerunning the DP unless the optional resegmentation scoring
mode is enabled.
\subsection{Prediction diagnostics}
\label{sec:prediction-diagnostics}

Because the predictive layer reuses segment-level posteriors, the group posterior
$P(g\mid\text{new})$ derived from \eqref{eq:group-lik-map} and the pointwise predictive summaries
inherit their calibration quality from the segmentation. Use two diagnostics for real-data
deployments:
\begin{itemize}\itemsep2pt
\item \textbf{Held-out log-likelihood (HLL) trace.} Split each test sequence into blocks and
compute the per-block posterior-predictive log-likelihood under \eqref{eq:pp-seg}. The sum over
blocks is the HLL for that sequence. Compare HLL trajectories across candidate segmentations
(exported MAP vs.\ Bayes curve) and across groups; a well-calibrated model yields an HLL
ranking consistent with the $k$-posterior ranking on held-out data.
\item \textbf{PIT residuals.} For point-valued new observations with continuous predictive CDF,
evaluate the probability-integral transform $u_i=F_g(y_i^{\text{new}}\mid\mathcal{M}_g)$.
Under correct calibration the $\{u_i\}$ are Uniform$(0,1)$; histograms or Kolmogorov--Smirnov
summaries flag miscalibration (overconfident predictive widths show up as U-shaped histograms).
\end{itemize}
These diagnostics are cheap post-hoc checks. In this paper they are planned outputs
for the real-data pipelines in Section~\ref{sec:experiments}; placeholder tables identify the planned
HLL and boundary-recovery fields that the finalized pipelines populate.
